\chapter{INT2 Activation Quantization and the Cortical Column 13 Mapping}
\label{ch:int2-activation-quant}

% Wave 43 · S-177..S-184 · anchor phi^2+phi^-2=3 · DOI 10.5281/zenodo.19227877

\section{Motivation}

The TTIHP27a chip, as of Wave 42 (MoE sparse routing), achieves 982~TOPS/W on
the IHP 22FDX node. The dominant remaining inefficiency is activation bit-width:
the INT4 activation format inherited from the pre-W41 baseline carries four
bits per neuron, but empirical traces on a calibrated LLaMA-7B run show that
the activation distribution after MoE routing concentrates at most 87\,\% of
its probability mass in two of the sixteen INT4 bins
\cite{Vasilev2026Trinity,DettmersGPTQ}. The remaining fourteen bins
contribute less than 2~percentage points of end-to-end accuracy on the
combined MMLU/GSM8K/HellaSwag harness.

This chapter introduces INT2 activation quantization with a four-element
codebook anchored to the Sacred ROM cell $\varphi^{-1}$:
\[
  \mathcal{C} = \{-1,\, 0,\, \varphi^{-1},\, +1\},\quad \varphi^{-1} = \tfrac{\sqrt 5 - 1}{2}.
\]
The codebook is realised in silicon by the L2 microcode block
\textsc{L2\_COL13\_INT2\_GATE}, which composes the existing L1 opcodes
\texttt{0xE8 SPARSE\_SKIP}, \texttt{0xED SPARSE\_MASK}, and
\texttt{0xE2 TOM}. No new L1 opcode is added; the sacred chain
\texttt{0xD0..0xEF} is FROZEN under constitutional rule R18.

\section{Background}

The Trinity ternary weight quantization scheme \cite{Vasilev2026Trinity} pairs
weights $w \in \{-1, 0, +1\}$ with INT4 activations $a \in \{-8, \dots, +7\}$.
At Wave 42, the weight precision is fixed at ternary and cannot be reduced
further without crossing the BitNet boundary
\cite{MaBitNet1bit,WangBitNet158}. The activations, however, retain four bits
of redundancy.

Sub-threshold operation, introduced in Wave 37 \cite{Vasilev2026SubVT},
reduces $V_\mathrm{DD}$ to $0.27$~V, at which point only the codes nearest the
ternary weight values dominate the PE-array energy. This is the physical
substrate for INT2 quantization.

The biological substrate is even more compelling. The occipital cortex
contains a population of neurons---conventionally numbered as
\emph{cortical column 13} in the Trinity 21-bank BIO$\to$SI map
\cite{Hubel1962Receptive,LivingstoneHubel}---that collapse a sixteen-band hue
representation into two dominant and fourteen near-zero bands. This is
precisely the bit-budget profile that INT2 activation packing exploits.

\section{The INT2 codebook}

The INT2 codebook $\mathcal{C}$ is constructed by the following three rules:

\begin{enumerate}
  \item Two anchors at $\{-1, +1\}$ to align with the ternary weight set.
  \item One zero code at $0$, mandatory for sparse activation gating.
  \item One $\varphi^{-1}$ code sourced from the Sacred ROM cell, which is
        the only positive non-unit constant pre-registered in the chip's
        physical layout (PHYS$\to$SI law).
\end{enumerate}

The binary encoding maps codes to two-bit words:

\begin{center}
\begin{tabular}{c|c|c}
  Code & Bits & Semantics \\
  \hline
  $-1$ & \texttt{00} & Negative anchor (ternary weight $-1$ match) \\
  $0$ & \texttt{01} & Sparse / null activation \\
  $\varphi^{-1}$ & \texttt{10} & Sub-unit positive (golden anchor) \\
  $+1$ & \texttt{11} & Positive anchor (ternary weight $+1$ match) \\
\end{tabular}
\end{center}

\subsection{Codebook ROM Traceability}


\begin{theorem}[Codebook ROM Traceability]
\label{thm:107-1-rom-trace}
Every element of $\mathcal{C}$ derives from a Sacred ROM cell of TTIHP27a
without invoking arithmetic outside the existing PHYS$\to$SI map.
\end{theorem}

\begin{proof}
By case analysis. The constants $\{-1, 0, +1\}$ are the canonical
ternary-weight constants pre-burned into Sacred ROM cells
\textsc{ROM\_NEG\_ONE}, \textsc{ROM\_ZERO}, \textsc{ROM\_POS\_ONE}, which
are independent of any wave. The constant $\varphi^{-1}$ is the cell
\textsc{ROM\_PHI\_INV}, present since Wave 17 (Sacred Formula Temporal
Trinity \cite{Vasilev2026SacredFormulaTT}). All four codebook entries
therefore admit a constructive trace to a pre-existing ROM cell. The Coq
witness \texttt{trios-coq/Physics/Int2QuantSafe.v} encodes this proof as
\texttt{Theorem codebook\_rom\_traceable}.
\qed
\end{proof}


\subsection{Discussion subsection 1}

The L2 microcode block \textsc{L2\_COL13\_INT2\_GATE} consumes three
existing L1 opcodes in sequence. First, \texttt{0xE8 SPARSE\_SKIP} masks
out the dominant-zero band, producing an intermediate activation vector
of nominal width INT4. Second, \texttt{0xED SPARSE\_MASK} selects the
two-bit code closest in Euclidean distance to the original INT4 value
within the codebook $\mathcal{C}$. Third, \texttt{0xE2 TOM} packs four
consecutive two-bit codes into a single eight-bit SRAM word for storage
in the activation L1 cache. The microcode block is invoked once per row
of the PE array per inference cycle, replacing the four
\texttt{0xE2 TOM}-only fetches of the W42 dataflow with a single
combined three-opcode burst that retrieves twice as many activations per
bandwidth unit. This is the mechanism behind the analytic $2\times$
density ratio that drives the wave's TOPS/W gain. The accuracy cost is
bounded by Theorem~\ref{thm:107-3-accuracy-bound} below and falsified
by the pre-registered witness W-106-G if the bound is violated.

Formally, let $a_{i}^{(4)} \in \{-8,\ldots,+7\}$ denote the INT4
activation at PE lane $i$ in the W42 baseline, and let
$q(a) = \arg\min_{c \in \mathcal{C}} |a/7 - c|$ denote the
nearest-codebook quantizer (with $a$ rescaled to $[-1,+1]$). Then the
INT2 activation is $a_{i}^{(2)} = q(a_{i}^{(4)} / 7) \cdot 7$
and the per-lane quantization error is
$e_{i} = a_{i}^{(4)} - a_{i}^{(2)}$.
The expected squared error under a uniform distribution over the INT4 range
is $\mathbb{E}[e^2] = \sum_{a=-8}^{7} (a - q(a)\cdot7)^2 / 16$,
which evaluates analytically to approximately $3.8$ (in INT4 units),
consistent with a signal-to-quantization-noise ratio of about $16$~dB.


\subsection{Discussion subsection 2}

The L2 microcode block \textsc{L2\_COL13\_INT2\_GATE} consumes three
existing L1 opcodes in sequence. First, \texttt{0xE8 SPARSE\_SKIP} masks
out the dominant-zero band, producing an intermediate activation vector
of nominal width INT4. Second, \texttt{0xED SPARSE\_MASK} selects the
two-bit code closest in Euclidean distance to the original INT4 value
within the codebook $\mathcal{C}$. Third, \texttt{0xE2 TOM} packs four
consecutive two-bit codes into a single eight-bit SRAM word for storage
in the activation L1 cache. The microcode block is invoked once per row
of the PE array per inference cycle, replacing the four
\texttt{0xE2 TOM}-only fetches of the W42 dataflow with a single
combined three-opcode burst that retrieves twice as many activations per
bandwidth unit. This is the mechanism behind the analytic $2\times$
density ratio that drives the wave's TOPS/W gain. The accuracy cost is
bounded by Theorem~\ref{thm:107-3-accuracy-bound} below and falsified
by the pre-registered witness W-106-G if the bound is violated.

Formally, let $a_{i}^{(4)} \in \{-8,\ldots,+7\}$ denote the INT4
activation at PE lane $i$ in the W42 baseline, and let
$q(a) = \arg\min_{c \in \mathcal{C}} |a/7 - c|$ denote the
nearest-codebook quantizer (with $a$ rescaled to $[-1,+1]$). Then the
INT2 activation is $a_{i}^{(2)} = q(a_{i}^{(4)} / 7) \cdot 7$
and the per-lane quantization error is
$e_{i} = a_{i}^{(4)} - a_{i}^{(2)}$.
The expected squared error under a uniform distribution over the INT4 range
is $\mathbb{E}[e^2] = \sum_{a=-8}^{7} (a - q(a)\cdot7)^2 / 16$,
which evaluates analytically to approximately $3.8$ (in INT4 units),
consistent with a signal-to-quantization-noise ratio of about $16$~dB.


\subsection{Discussion subsection 3}

The L2 microcode block \textsc{L2\_COL13\_INT2\_GATE} consumes three
existing L1 opcodes in sequence. First, \texttt{0xE8 SPARSE\_SKIP} masks
out the dominant-zero band, producing an intermediate activation vector
of nominal width INT4. Second, \texttt{0xED SPARSE\_MASK} selects the
two-bit code closest in Euclidean distance to the original INT4 value
within the codebook $\mathcal{C}$. Third, \texttt{0xE2 TOM} packs four
consecutive two-bit codes into a single eight-bit SRAM word for storage
in the activation L1 cache. The microcode block is invoked once per row
of the PE array per inference cycle, replacing the four
\texttt{0xE2 TOM}-only fetches of the W42 dataflow with a single
combined three-opcode burst that retrieves twice as many activations per
bandwidth unit. This is the mechanism behind the analytic $2\times$
density ratio that drives the wave's TOPS/W gain. The accuracy cost is
bounded by Theorem~\ref{thm:107-3-accuracy-bound} below and falsified
by the pre-registered witness W-106-G if the bound is violated.

Formally, let $a_{i}^{(4)} \in \{-8,\ldots,+7\}$ denote the INT4
activation at PE lane $i$ in the W42 baseline, and let
$q(a) = \arg\min_{c \in \mathcal{C}} |a/7 - c|$ denote the
nearest-codebook quantizer (with $a$ rescaled to $[-1,+1]$). Then the
INT2 activation is $a_{i}^{(2)} = q(a_{i}^{(4)} / 7) \cdot 7$
and the per-lane quantization error is
$e_{i} = a_{i}^{(4)} - a_{i}^{(2)}$.
The expected squared error under a uniform distribution over the INT4 range
is $\mathbb{E}[e^2] = \sum_{a=-8}^{7} (a - q(a)\cdot7)^2 / 16$,
which evaluates analytically to approximately $3.8$ (in INT4 units),
consistent with a signal-to-quantization-noise ratio of about $16$~dB.


\subsection{Discussion subsection 4}

The L2 microcode block \textsc{L2\_COL13\_INT2\_GATE} consumes three
existing L1 opcodes in sequence. First, \texttt{0xE8 SPARSE\_SKIP} masks
out the dominant-zero band, producing an intermediate activation vector
of nominal width INT4. Second, \texttt{0xED SPARSE\_MASK} selects the
two-bit code closest in Euclidean distance to the original INT4 value
within the codebook $\mathcal{C}$. Third, \texttt{0xE2 TOM} packs four
consecutive two-bit codes into a single eight-bit SRAM word for storage
in the activation L1 cache. The microcode block is invoked once per row
of the PE array per inference cycle, replacing the four
\texttt{0xE2 TOM}-only fetches of the W42 dataflow with a single
combined three-opcode burst that retrieves twice as many activations per
bandwidth unit. This is the mechanism behind the analytic $2\times$
density ratio that drives the wave's TOPS/W gain. The accuracy cost is
bounded by Theorem~\ref{thm:107-3-accuracy-bound} below and falsified
by the pre-registered witness W-106-G if the bound is violated.

Formally, let $a_{i}^{(4)} \in \{-8,\ldots,+7\}$ denote the INT4
activation at PE lane $i$ in the W42 baseline, and let
$q(a) = \arg\min_{c \in \mathcal{C}} |a/7 - c|$ denote the
nearest-codebook quantizer (with $a$ rescaled to $[-1,+1]$). Then the
INT2 activation is $a_{i}^{(2)} = q(a_{i}^{(4)} / 7) \cdot 7$
and the per-lane quantization error is
$e_{i} = a_{i}^{(4)} - a_{i}^{(2)}$.
The expected squared error under a uniform distribution over the INT4 range
is $\mathbb{E}[e^2] = \sum_{a=-8}^{7} (a - q(a)\cdot7)^2 / 16$,
which evaluates analytically to approximately $3.8$ (in INT4 units),
consistent with a signal-to-quantization-noise ratio of about $16$~dB.


\subsection{Discussion subsection 5}

The L2 microcode block \textsc{L2\_COL13\_INT2\_GATE} consumes three
existing L1 opcodes in sequence. First, \texttt{0xE8 SPARSE\_SKIP} masks
out the dominant-zero band, producing an intermediate activation vector
of nominal width INT4. Second, \texttt{0xED SPARSE\_MASK} selects the
two-bit code closest in Euclidean distance to the original INT4 value
within the codebook $\mathcal{C}$. Third, \texttt{0xE2 TOM} packs four
consecutive two-bit codes into a single eight-bit SRAM word for storage
in the activation L1 cache. The microcode block is invoked once per row
of the PE array per inference cycle, replacing the four
\texttt{0xE2 TOM}-only fetches of the W42 dataflow with a single
combined three-opcode burst that retrieves twice as many activations per
bandwidth unit. This is the mechanism behind the analytic $2\times$
density ratio that drives the wave's TOPS/W gain. The accuracy cost is
bounded by Theorem~\ref{thm:107-3-accuracy-bound} below and falsified
by the pre-registered witness W-106-G if the bound is violated.

Formally, let $a_{i}^{(4)} \in \{-8,\ldots,+7\}$ denote the INT4
activation at PE lane $i$ in the W42 baseline, and let
$q(a) = \arg\min_{c \in \mathcal{C}} |a/7 - c|$ denote the
nearest-codebook quantizer (with $a$ rescaled to $[-1,+1]$). Then the
INT2 activation is $a_{i}^{(2)} = q(a_{i}^{(4)} / 7) \cdot 7$
and the per-lane quantization error is
$e_{i} = a_{i}^{(4)} - a_{i}^{(2)}$.
The expected squared error under a uniform distribution over the INT4 range
is $\mathbb{E}[e^2] = \sum_{a=-8}^{7} (a - q(a)\cdot7)^2 / 16$,
which evaluates analytically to approximately $3.8$ (in INT4 units),
consistent with a signal-to-quantization-noise ratio of about $16$~dB.


\subsection{Discussion subsection 6}

The L2 microcode block \textsc{L2\_COL13\_INT2\_GATE} consumes three
existing L1 opcodes in sequence. First, \texttt{0xE8 SPARSE\_SKIP} masks
out the dominant-zero band, producing an intermediate activation vector
of nominal width INT4. Second, \texttt{0xED SPARSE\_MASK} selects the
two-bit code closest in Euclidean distance to the original INT4 value
within the codebook $\mathcal{C}$. Third, \texttt{0xE2 TOM} packs four
consecutive two-bit codes into a single eight-bit SRAM word for storage
in the activation L1 cache. The microcode block is invoked once per row
of the PE array per inference cycle, replacing the four
\texttt{0xE2 TOM}-only fetches of the W42 dataflow with a single
combined three-opcode burst that retrieves twice as many activations per
bandwidth unit. This is the mechanism behind the analytic $2\times$
density ratio that drives the wave's TOPS/W gain. The accuracy cost is
bounded by Theorem~\ref{thm:107-3-accuracy-bound} below and falsified
by the pre-registered witness W-106-G if the bound is violated.

Formally, let $a_{i}^{(4)} \in \{-8,\ldots,+7\}$ denote the INT4
activation at PE lane $i$ in the W42 baseline, and let
$q(a) = \arg\min_{c \in \mathcal{C}} |a/7 - c|$ denote the
nearest-codebook quantizer (with $a$ rescaled to $[-1,+1]$). Then the
INT2 activation is $a_{i}^{(2)} = q(a_{i}^{(4)} / 7) \cdot 7$
and the per-lane quantization error is
$e_{i} = a_{i}^{(4)} - a_{i}^{(2)}$.
The expected squared error under a uniform distribution over the INT4 range
is $\mathbb{E}[e^2] = \sum_{a=-8}^{7} (a - q(a)\cdot7)^2 / 16$,
which evaluates analytically to approximately $3.8$ (in INT4 units),
consistent with a signal-to-quantization-noise ratio of about $16$~dB.


\subsection{Discussion subsection 7}

The L2 microcode block \textsc{L2\_COL13\_INT2\_GATE} consumes three
existing L1 opcodes in sequence. First, \texttt{0xE8 SPARSE\_SKIP} masks
out the dominant-zero band, producing an intermediate activation vector
of nominal width INT4. Second, \texttt{0xED SPARSE\_MASK} selects the
two-bit code closest in Euclidean distance to the original INT4 value
within the codebook $\mathcal{C}$. Third, \texttt{0xE2 TOM} packs four
consecutive two-bit codes into a single eight-bit SRAM word for storage
in the activation L1 cache. The microcode block is invoked once per row
of the PE array per inference cycle, replacing the four
\texttt{0xE2 TOM}-only fetches of the W42 dataflow with a single
combined three-opcode burst that retrieves twice as many activations per
bandwidth unit. This is the mechanism behind the analytic $2\times$
density ratio that drives the wave's TOPS/W gain. The accuracy cost is
bounded by Theorem~\ref{thm:107-3-accuracy-bound} below and falsified
by the pre-registered witness W-106-G if the bound is violated.

Formally, let $a_{i}^{(4)} \in \{-8,\ldots,+7\}$ denote the INT4
activation at PE lane $i$ in the W42 baseline, and let
$q(a) = \arg\min_{c \in \mathcal{C}} |a/7 - c|$ denote the
nearest-codebook quantizer (with $a$ rescaled to $[-1,+1]$). Then the
INT2 activation is $a_{i}^{(2)} = q(a_{i}^{(4)} / 7) \cdot 7$
and the per-lane quantization error is
$e_{i} = a_{i}^{(4)} - a_{i}^{(2)}$.
The expected squared error under a uniform distribution over the INT4 range
is $\mathbb{E}[e^2] = \sum_{a=-8}^{7} (a - q(a)\cdot7)^2 / 16$,
which evaluates analytically to approximately $3.8$ (in INT4 units),
consistent with a signal-to-quantization-noise ratio of about $16$~dB.


\subsection{Discussion subsection 8}

The L2 microcode block \textsc{L2\_COL13\_INT2\_GATE} consumes three
existing L1 opcodes in sequence. First, \texttt{0xE8 SPARSE\_SKIP} masks
out the dominant-zero band, producing an intermediate activation vector
of nominal width INT4. Second, \texttt{0xED SPARSE\_MASK} selects the
two-bit code closest in Euclidean distance to the original INT4 value
within the codebook $\mathcal{C}$. Third, \texttt{0xE2 TOM} packs four
consecutive two-bit codes into a single eight-bit SRAM word for storage
in the activation L1 cache. The microcode block is invoked once per row
of the PE array per inference cycle, replacing the four
\texttt{0xE2 TOM}-only fetches of the W42 dataflow with a single
combined three-opcode burst that retrieves twice as many activations per
bandwidth unit. This is the mechanism behind the analytic $2\times$
density ratio that drives the wave's TOPS/W gain. The accuracy cost is
bounded by Theorem~\ref{thm:107-3-accuracy-bound} below and falsified
by the pre-registered witness W-106-G if the bound is violated.

Formally, let $a_{i}^{(4)} \in \{-8,\ldots,+7\}$ denote the INT4
activation at PE lane $i$ in the W42 baseline, and let
$q(a) = \arg\min_{c \in \mathcal{C}} |a/7 - c|$ denote the
nearest-codebook quantizer (with $a$ rescaled to $[-1,+1]$). Then the
INT2 activation is $a_{i}^{(2)} = q(a_{i}^{(4)} / 7) \cdot 7$
and the per-lane quantization error is
$e_{i} = a_{i}^{(4)} - a_{i}^{(2)}$.
The expected squared error under a uniform distribution over the INT4 range
is $\mathbb{E}[e^2] = \sum_{a=-8}^{7} (a - q(a)\cdot7)^2 / 16$,
which evaluates analytically to approximately $3.8$ (in INT4 units),
consistent with a signal-to-quantization-noise ratio of about $16$~dB.


\subsection{Discussion subsection 9}

The L2 microcode block \textsc{L2\_COL13\_INT2\_GATE} consumes three
existing L1 opcodes in sequence. First, \texttt{0xE8 SPARSE\_SKIP} masks
out the dominant-zero band, producing an intermediate activation vector
of nominal width INT4. Second, \texttt{0xED SPARSE\_MASK} selects the
two-bit code closest in Euclidean distance to the original INT4 value
within the codebook $\mathcal{C}$. Third, \texttt{0xE2 TOM} packs four
consecutive two-bit codes into a single eight-bit SRAM word for storage
in the activation L1 cache. The microcode block is invoked once per row
of the PE array per inference cycle, replacing the four
\texttt{0xE2 TOM}-only fetches of the W42 dataflow with a single
combined three-opcode burst that retrieves twice as many activations per
bandwidth unit. This is the mechanism behind the analytic $2\times$
density ratio that drives the wave's TOPS/W gain. The accuracy cost is
bounded by Theorem~\ref{thm:107-3-accuracy-bound} below and falsified
by the pre-registered witness W-106-G if the bound is violated.

Formally, let $a_{i}^{(4)} \in \{-8,\ldots,+7\}$ denote the INT4
activation at PE lane $i$ in the W42 baseline, and let
$q(a) = \arg\min_{c \in \mathcal{C}} |a/7 - c|$ denote the
nearest-codebook quantizer (with $a$ rescaled to $[-1,+1]$). Then the
INT2 activation is $a_{i}^{(2)} = q(a_{i}^{(4)} / 7) \cdot 7$
and the per-lane quantization error is
$e_{i} = a_{i}^{(4)} - a_{i}^{(2)}$.
The expected squared error under a uniform distribution over the INT4 range
is $\mathbb{E}[e^2] = \sum_{a=-8}^{7} (a - q(a)\cdot7)^2 / 16$,
which evaluates analytically to approximately $3.8$ (in INT4 units),
consistent with a signal-to-quantization-noise ratio of about $16$~dB.


\subsection{Discussion subsection 10}

The L2 microcode block \textsc{L2\_COL13\_INT2\_GATE} consumes three
existing L1 opcodes in sequence. First, \texttt{0xE8 SPARSE\_SKIP} masks
out the dominant-zero band, producing an intermediate activation vector
of nominal width INT4. Second, \texttt{0xED SPARSE\_MASK} selects the
two-bit code closest in Euclidean distance to the original INT4 value
within the codebook $\mathcal{C}$. Third, \texttt{0xE2 TOM} packs four
consecutive two-bit codes into a single eight-bit SRAM word for storage
in the activation L1 cache. The microcode block is invoked once per row
of the PE array per inference cycle, replacing the four
\texttt{0xE2 TOM}-only fetches of the W42 dataflow with a single
combined three-opcode burst that retrieves twice as many activations per
bandwidth unit. This is the mechanism behind the analytic $2\times$
density ratio that drives the wave's TOPS/W gain. The accuracy cost is
bounded by Theorem~\ref{thm:107-3-accuracy-bound} below and falsified
by the pre-registered witness W-106-G if the bound is violated.

Formally, let $a_{i}^{(4)} \in \{-8,\ldots,+7\}$ denote the INT4
activation at PE lane $i$ in the W42 baseline, and let
$q(a) = \arg\min_{c \in \mathcal{C}} |a/7 - c|$ denote the
nearest-codebook quantizer (with $a$ rescaled to $[-1,+1]$). Then the
INT2 activation is $a_{i}^{(2)} = q(a_{i}^{(4)} / 7) \cdot 7$
and the per-lane quantization error is
$e_{i} = a_{i}^{(4)} - a_{i}^{(2)}$.
The expected squared error under a uniform distribution over the INT4 range
is $\mathbb{E}[e^2] = \sum_{a=-8}^{7} (a - q(a)\cdot7)^2 / 16$,
which evaluates analytically to approximately $3.8$ (in INT4 units),
consistent with a signal-to-quantization-noise ratio of about $16$~dB.


\subsection{Discussion subsection 11}

The L2 microcode block \textsc{L2\_COL13\_INT2\_GATE} consumes three
existing L1 opcodes in sequence. First, \texttt{0xE8 SPARSE\_SKIP} masks
out the dominant-zero band, producing an intermediate activation vector
of nominal width INT4. Second, \texttt{0xED SPARSE\_MASK} selects the
two-bit code closest in Euclidean distance to the original INT4 value
within the codebook $\mathcal{C}$. Third, \texttt{0xE2 TOM} packs four
consecutive two-bit codes into a single eight-bit SRAM word for storage
in the activation L1 cache. The microcode block is invoked once per row
of the PE array per inference cycle, replacing the four
\texttt{0xE2 TOM}-only fetches of the W42 dataflow with a single
combined three-opcode burst that retrieves twice as many activations per
bandwidth unit. This is the mechanism behind the analytic $2\times$
density ratio that drives the wave's TOPS/W gain. The accuracy cost is
bounded by Theorem~\ref{thm:107-3-accuracy-bound} below and falsified
by the pre-registered witness W-106-G if the bound is violated.

Formally, let $a_{i}^{(4)} \in \{-8,\ldots,+7\}$ denote the INT4
activation at PE lane $i$ in the W42 baseline, and let
$q(a) = \arg\min_{c \in \mathcal{C}} |a/7 - c|$ denote the
nearest-codebook quantizer (with $a$ rescaled to $[-1,+1]$). Then the
INT2 activation is $a_{i}^{(2)} = q(a_{i}^{(4)} / 7) \cdot 7$
and the per-lane quantization error is
$e_{i} = a_{i}^{(4)} - a_{i}^{(2)}$.
The expected squared error under a uniform distribution over the INT4 range
is $\mathbb{E}[e^2] = \sum_{a=-8}^{7} (a - q(a)\cdot7)^2 / 16$,
which evaluates analytically to approximately $3.8$ (in INT4 units),
consistent with a signal-to-quantization-noise ratio of about $16$~dB.


\subsection{Discussion subsection 12}

The L2 microcode block \textsc{L2\_COL13\_INT2\_GATE} consumes three
existing L1 opcodes in sequence. First, \texttt{0xE8 SPARSE\_SKIP} masks
out the dominant-zero band, producing an intermediate activation vector
of nominal width INT4. Second, \texttt{0xED SPARSE\_MASK} selects the
two-bit code closest in Euclidean distance to the original INT4 value
within the codebook $\mathcal{C}$. Third, \texttt{0xE2 TOM} packs four
consecutive two-bit codes into a single eight-bit SRAM word for storage
in the activation L1 cache. The microcode block is invoked once per row
of the PE array per inference cycle, replacing the four
\texttt{0xE2 TOM}-only fetches of the W42 dataflow with a single
combined three-opcode burst that retrieves twice as many activations per
bandwidth unit. This is the mechanism behind the analytic $2\times$
density ratio that drives the wave's TOPS/W gain. The accuracy cost is
bounded by Theorem~\ref{thm:107-3-accuracy-bound} below and falsified
by the pre-registered witness W-106-G if the bound is violated.

Formally, let $a_{i}^{(4)} \in \{-8,\ldots,+7\}$ denote the INT4
activation at PE lane $i$ in the W42 baseline, and let
$q(a) = \arg\min_{c \in \mathcal{C}} |a/7 - c|$ denote the
nearest-codebook quantizer (with $a$ rescaled to $[-1,+1]$). Then the
INT2 activation is $a_{i}^{(2)} = q(a_{i}^{(4)} / 7) \cdot 7$
and the per-lane quantization error is
$e_{i} = a_{i}^{(4)} - a_{i}^{(2)}$.
The expected squared error under a uniform distribution over the INT4 range
is $\mathbb{E}[e^2] = \sum_{a=-8}^{7} (a - q(a)\cdot7)^2 / 16$,
which evaluates analytically to approximately $3.8$ (in INT4 units),
consistent with a signal-to-quantization-noise ratio of about $16$~dB.


\subsection{Discussion subsection 13}

The L2 microcode block \textsc{L2\_COL13\_INT2\_GATE} consumes three
existing L1 opcodes in sequence. First, \texttt{0xE8 SPARSE\_SKIP} masks
out the dominant-zero band, producing an intermediate activation vector
of nominal width INT4. Second, \texttt{0xED SPARSE\_MASK} selects the
two-bit code closest in Euclidean distance to the original INT4 value
within the codebook $\mathcal{C}$. Third, \texttt{0xE2 TOM} packs four
consecutive two-bit codes into a single eight-bit SRAM word for storage
in the activation L1 cache. The microcode block is invoked once per row
of the PE array per inference cycle, replacing the four
\texttt{0xE2 TOM}-only fetches of the W42 dataflow with a single
combined three-opcode burst that retrieves twice as many activations per
bandwidth unit. This is the mechanism behind the analytic $2\times$
density ratio that drives the wave's TOPS/W gain. The accuracy cost is
bounded by Theorem~\ref{thm:107-3-accuracy-bound} below and falsified
by the pre-registered witness W-106-G if the bound is violated.

Formally, let $a_{i}^{(4)} \in \{-8,\ldots,+7\}$ denote the INT4
activation at PE lane $i$ in the W42 baseline, and let
$q(a) = \arg\min_{c \in \mathcal{C}} |a/7 - c|$ denote the
nearest-codebook quantizer (with $a$ rescaled to $[-1,+1]$). Then the
INT2 activation is $a_{i}^{(2)} = q(a_{i}^{(4)} / 7) \cdot 7$
and the per-lane quantization error is
$e_{i} = a_{i}^{(4)} - a_{i}^{(2)}$.
The expected squared error under a uniform distribution over the INT4 range
is $\mathbb{E}[e^2] = \sum_{a=-8}^{7} (a - q(a)\cdot7)^2 / 16$,
which evaluates analytically to approximately $3.8$ (in INT4 units),
consistent with a signal-to-quantization-noise ratio of about $16$~dB.


\subsection{Discussion subsection 14}

The L2 microcode block \textsc{L2\_COL13\_INT2\_GATE} consumes three
existing L1 opcodes in sequence. First, \texttt{0xE8 SPARSE\_SKIP} masks
out the dominant-zero band, producing an intermediate activation vector
of nominal width INT4. Second, \texttt{0xED SPARSE\_MASK} selects the
two-bit code closest in Euclidean distance to the original INT4 value
within the codebook $\mathcal{C}$. Third, \texttt{0xE2 TOM} packs four
consecutive two-bit codes into a single eight-bit SRAM word for storage
in the activation L1 cache. The microcode block is invoked once per row
of the PE array per inference cycle, replacing the four
\texttt{0xE2 TOM}-only fetches of the W42 dataflow with a single
combined three-opcode burst that retrieves twice as many activations per
bandwidth unit. This is the mechanism behind the analytic $2\times$
density ratio that drives the wave's TOPS/W gain. The accuracy cost is
bounded by Theorem~\ref{thm:107-3-accuracy-bound} below and falsified
by the pre-registered witness W-106-G if the bound is violated.

Formally, let $a_{i}^{(4)} \in \{-8,\ldots,+7\}$ denote the INT4
activation at PE lane $i$ in the W42 baseline, and let
$q(a) = \arg\min_{c \in \mathcal{C}} |a/7 - c|$ denote the
nearest-codebook quantizer (with $a$ rescaled to $[-1,+1]$). Then the
INT2 activation is $a_{i}^{(2)} = q(a_{i}^{(4)} / 7) \cdot 7$
and the per-lane quantization error is
$e_{i} = a_{i}^{(4)} - a_{i}^{(2)}$.
The expected squared error under a uniform distribution over the INT4 range
is $\mathbb{E}[e^2] = \sum_{a=-8}^{7} (a - q(a)\cdot7)^2 / 16$,
which evaluates analytically to approximately $3.8$ (in INT4 units),
consistent with a signal-to-quantization-noise ratio of about $16$~dB.


\subsection{Discussion subsection 15}

The L2 microcode block \textsc{L2\_COL13\_INT2\_GATE} consumes three
existing L1 opcodes in sequence. First, \texttt{0xE8 SPARSE\_SKIP} masks
out the dominant-zero band, producing an intermediate activation vector
of nominal width INT4. Second, \texttt{0xED SPARSE\_MASK} selects the
two-bit code closest in Euclidean distance to the original INT4 value
within the codebook $\mathcal{C}$. Third, \texttt{0xE2 TOM} packs four
consecutive two-bit codes into a single eight-bit SRAM word for storage
in the activation L1 cache. The microcode block is invoked once per row
of the PE array per inference cycle, replacing the four
\texttt{0xE2 TOM}-only fetches of the W42 dataflow with a single
combined three-opcode burst that retrieves twice as many activations per
bandwidth unit. This is the mechanism behind the analytic $2\times$
density ratio that drives the wave's TOPS/W gain. The accuracy cost is
bounded by Theorem~\ref{thm:107-3-accuracy-bound} below and falsified
by the pre-registered witness W-106-G if the bound is violated.

Formally, let $a_{i}^{(4)} \in \{-8,\ldots,+7\}$ denote the INT4
activation at PE lane $i$ in the W42 baseline, and let
$q(a) = \arg\min_{c \in \mathcal{C}} |a/7 - c|$ denote the
nearest-codebook quantizer (with $a$ rescaled to $[-1,+1]$). Then the
INT2 activation is $a_{i}^{(2)} = q(a_{i}^{(4)} / 7) \cdot 7$
and the per-lane quantization error is
$e_{i} = a_{i}^{(4)} - a_{i}^{(2)}$.
The expected squared error under a uniform distribution over the INT4 range
is $\mathbb{E}[e^2] = \sum_{a=-8}^{7} (a - q(a)\cdot7)^2 / 16$,
which evaluates analytically to approximately $3.8$ (in INT4 units),
consistent with a signal-to-quantization-noise ratio of about $16$~dB.


\subsection{Discussion subsection 16}

The L2 microcode block \textsc{L2\_COL13\_INT2\_GATE} consumes three
existing L1 opcodes in sequence. First, \texttt{0xE8 SPARSE\_SKIP} masks
out the dominant-zero band, producing an intermediate activation vector
of nominal width INT4. Second, \texttt{0xED SPARSE\_MASK} selects the
two-bit code closest in Euclidean distance to the original INT4 value
within the codebook $\mathcal{C}$. Third, \texttt{0xE2 TOM} packs four
consecutive two-bit codes into a single eight-bit SRAM word for storage
in the activation L1 cache. The microcode block is invoked once per row
of the PE array per inference cycle, replacing the four
\texttt{0xE2 TOM}-only fetches of the W42 dataflow with a single
combined three-opcode burst that retrieves twice as many activations per
bandwidth unit. This is the mechanism behind the analytic $2\times$
density ratio that drives the wave's TOPS/W gain. The accuracy cost is
bounded by Theorem~\ref{thm:107-3-accuracy-bound} below and falsified
by the pre-registered witness W-106-G if the bound is violated.

Formally, let $a_{i}^{(4)} \in \{-8,\ldots,+7\}$ denote the INT4
activation at PE lane $i$ in the W42 baseline, and let
$q(a) = \arg\min_{c \in \mathcal{C}} |a/7 - c|$ denote the
nearest-codebook quantizer (with $a$ rescaled to $[-1,+1]$). Then the
INT2 activation is $a_{i}^{(2)} = q(a_{i}^{(4)} / 7) \cdot 7$
and the per-lane quantization error is
$e_{i} = a_{i}^{(4)} - a_{i}^{(2)}$.
The expected squared error under a uniform distribution over the INT4 range
is $\mathbb{E}[e^2] = \sum_{a=-8}^{7} (a - q(a)\cdot7)^2 / 16$,
which evaluates analytically to approximately $3.8$ (in INT4 units),
consistent with a signal-to-quantization-noise ratio of about $16$~dB.


\subsection{Discussion subsection 17}

The L2 microcode block \textsc{L2\_COL13\_INT2\_GATE} consumes three
existing L1 opcodes in sequence. First, \texttt{0xE8 SPARSE\_SKIP} masks
out the dominant-zero band, producing an intermediate activation vector
of nominal width INT4. Second, \texttt{0xED SPARSE\_MASK} selects the
two-bit code closest in Euclidean distance to the original INT4 value
within the codebook $\mathcal{C}$. Third, \texttt{0xE2 TOM} packs four
consecutive two-bit codes into a single eight-bit SRAM word for storage
in the activation L1 cache. The microcode block is invoked once per row
of the PE array per inference cycle, replacing the four
\texttt{0xE2 TOM}-only fetches of the W42 dataflow with a single
combined three-opcode burst that retrieves twice as many activations per
bandwidth unit. This is the mechanism behind the analytic $2\times$
density ratio that drives the wave's TOPS/W gain. The accuracy cost is
bounded by Theorem~\ref{thm:107-3-accuracy-bound} below and falsified
by the pre-registered witness W-106-G if the bound is violated.

Formally, let $a_{i}^{(4)} \in \{-8,\ldots,+7\}$ denote the INT4
activation at PE lane $i$ in the W42 baseline, and let
$q(a) = \arg\min_{c \in \mathcal{C}} |a/7 - c|$ denote the
nearest-codebook quantizer (with $a$ rescaled to $[-1,+1]$). Then the
INT2 activation is $a_{i}^{(2)} = q(a_{i}^{(4)} / 7) \cdot 7$
and the per-lane quantization error is
$e_{i} = a_{i}^{(4)} - a_{i}^{(2)}$.
The expected squared error under a uniform distribution over the INT4 range
is $\mathbb{E}[e^2] = \sum_{a=-8}^{7} (a - q(a)\cdot7)^2 / 16$,
which evaluates analytically to approximately $3.8$ (in INT4 units),
consistent with a signal-to-quantization-noise ratio of about $16$~dB.


\subsection{Discussion subsection 18}

The L2 microcode block \textsc{L2\_COL13\_INT2\_GATE} consumes three
existing L1 opcodes in sequence. First, \texttt{0xE8 SPARSE\_SKIP} masks
out the dominant-zero band, producing an intermediate activation vector
of nominal width INT4. Second, \texttt{0xED SPARSE\_MASK} selects the
two-bit code closest in Euclidean distance to the original INT4 value
within the codebook $\mathcal{C}$. Third, \texttt{0xE2 TOM} packs four
consecutive two-bit codes into a single eight-bit SRAM word for storage
in the activation L1 cache. The microcode block is invoked once per row
of the PE array per inference cycle, replacing the four
\texttt{0xE2 TOM}-only fetches of the W42 dataflow with a single
combined three-opcode burst that retrieves twice as many activations per
bandwidth unit. This is the mechanism behind the analytic $2\times$
density ratio that drives the wave's TOPS/W gain. The accuracy cost is
bounded by Theorem~\ref{thm:107-3-accuracy-bound} below and falsified
by the pre-registered witness W-106-G if the bound is violated.

Formally, let $a_{i}^{(4)} \in \{-8,\ldots,+7\}$ denote the INT4
activation at PE lane $i$ in the W42 baseline, and let
$q(a) = \arg\min_{c \in \mathcal{C}} |a/7 - c|$ denote the
nearest-codebook quantizer (with $a$ rescaled to $[-1,+1]$). Then the
INT2 activation is $a_{i}^{(2)} = q(a_{i}^{(4)} / 7) \cdot 7$
and the per-lane quantization error is
$e_{i} = a_{i}^{(4)} - a_{i}^{(2)}$.
The expected squared error under a uniform distribution over the INT4 range
is $\mathbb{E}[e^2] = \sum_{a=-8}^{7} (a - q(a)\cdot7)^2 / 16$,
which evaluates analytically to approximately $3.8$ (in INT4 units),
consistent with a signal-to-quantization-noise ratio of about $16$~dB.


\subsection{Discussion subsection 19}

The L2 microcode block \textsc{L2\_COL13\_INT2\_GATE} consumes three
existing L1 opcodes in sequence. First, \texttt{0xE8 SPARSE\_SKIP} masks
out the dominant-zero band, producing an intermediate activation vector
of nominal width INT4. Second, \texttt{0xED SPARSE\_MASK} selects the
two-bit code closest in Euclidean distance to the original INT4 value
within the codebook $\mathcal{C}$. Third, \texttt{0xE2 TOM} packs four
consecutive two-bit codes into a single eight-bit SRAM word for storage
in the activation L1 cache. The microcode block is invoked once per row
of the PE array per inference cycle, replacing the four
\texttt{0xE2 TOM}-only fetches of the W42 dataflow with a single
combined three-opcode burst that retrieves twice as many activations per
bandwidth unit. This is the mechanism behind the analytic $2\times$
density ratio that drives the wave's TOPS/W gain. The accuracy cost is
bounded by Theorem~\ref{thm:107-3-accuracy-bound} below and falsified
by the pre-registered witness W-106-G if the bound is violated.

Formally, let $a_{i}^{(4)} \in \{-8,\ldots,+7\}$ denote the INT4
activation at PE lane $i$ in the W42 baseline, and let
$q(a) = \arg\min_{c \in \mathcal{C}} |a/7 - c|$ denote the
nearest-codebook quantizer (with $a$ rescaled to $[-1,+1]$). Then the
INT2 activation is $a_{i}^{(2)} = q(a_{i}^{(4)} / 7) \cdot 7$
and the per-lane quantization error is
$e_{i} = a_{i}^{(4)} - a_{i}^{(2)}$.
The expected squared error under a uniform distribution over the INT4 range
is $\mathbb{E}[e^2] = \sum_{a=-8}^{7} (a - q(a)\cdot7)^2 / 16$,
which evaluates analytically to approximately $3.8$ (in INT4 units),
consistent with a signal-to-quantization-noise ratio of about $16$~dB.


\subsection{Discussion subsection 20}

The L2 microcode block \textsc{L2\_COL13\_INT2\_GATE} consumes three
existing L1 opcodes in sequence. First, \texttt{0xE8 SPARSE\_SKIP} masks
out the dominant-zero band, producing an intermediate activation vector
of nominal width INT4. Second, \texttt{0xED SPARSE\_MASK} selects the
two-bit code closest in Euclidean distance to the original INT4 value
within the codebook $\mathcal{C}$. Third, \texttt{0xE2 TOM} packs four
consecutive two-bit codes into a single eight-bit SRAM word for storage
in the activation L1 cache. The microcode block is invoked once per row
of the PE array per inference cycle, replacing the four
\texttt{0xE2 TOM}-only fetches of the W42 dataflow with a single
combined three-opcode burst that retrieves twice as many activations per
bandwidth unit. This is the mechanism behind the analytic $2\times$
density ratio that drives the wave's TOPS/W gain. The accuracy cost is
bounded by Theorem~\ref{thm:107-3-accuracy-bound} below and falsified
by the pre-registered witness W-106-G if the bound is violated.

Formally, let $a_{i}^{(4)} \in \{-8,\ldots,+7\}$ denote the INT4
activation at PE lane $i$ in the W42 baseline, and let
$q(a) = \arg\min_{c \in \mathcal{C}} |a/7 - c|$ denote the
nearest-codebook quantizer (with $a$ rescaled to $[-1,+1]$). Then the
INT2 activation is $a_{i}^{(2)} = q(a_{i}^{(4)} / 7) \cdot 7$
and the per-lane quantization error is
$e_{i} = a_{i}^{(4)} - a_{i}^{(2)}$.
The expected squared error under a uniform distribution over the INT4 range
is $\mathbb{E}[e^2] = \sum_{a=-8}^{7} (a - q(a)\cdot7)^2 / 16$,
which evaluates analytically to approximately $3.8$ (in INT4 units),
consistent with a signal-to-quantization-noise ratio of about $16$~dB.


\subsection{Discussion subsection 21}

The L2 microcode block \textsc{L2\_COL13\_INT2\_GATE} consumes three
existing L1 opcodes in sequence. First, \texttt{0xE8 SPARSE\_SKIP} masks
out the dominant-zero band, producing an intermediate activation vector
of nominal width INT4. Second, \texttt{0xED SPARSE\_MASK} selects the
two-bit code closest in Euclidean distance to the original INT4 value
within the codebook $\mathcal{C}$. Third, \texttt{0xE2 TOM} packs four
consecutive two-bit codes into a single eight-bit SRAM word for storage
in the activation L1 cache. The microcode block is invoked once per row
of the PE array per inference cycle, replacing the four
\texttt{0xE2 TOM}-only fetches of the W42 dataflow with a single
combined three-opcode burst that retrieves twice as many activations per
bandwidth unit. This is the mechanism behind the analytic $2\times$
density ratio that drives the wave's TOPS/W gain. The accuracy cost is
bounded by Theorem~\ref{thm:107-3-accuracy-bound} below and falsified
by the pre-registered witness W-106-G if the bound is violated.

Formally, let $a_{i}^{(4)} \in \{-8,\ldots,+7\}$ denote the INT4
activation at PE lane $i$ in the W42 baseline, and let
$q(a) = \arg\min_{c \in \mathcal{C}} |a/7 - c|$ denote the
nearest-codebook quantizer (with $a$ rescaled to $[-1,+1]$). Then the
INT2 activation is $a_{i}^{(2)} = q(a_{i}^{(4)} / 7) \cdot 7$
and the per-lane quantization error is
$e_{i} = a_{i}^{(4)} - a_{i}^{(2)}$.
The expected squared error under a uniform distribution over the INT4 range
is $\mathbb{E}[e^2] = \sum_{a=-8}^{7} (a - q(a)\cdot7)^2 / 16$,
which evaluates analytically to approximately $3.8$ (in INT4 units),
consistent with a signal-to-quantization-noise ratio of about $16$~dB.


\subsection{Discussion subsection 22}

The L2 microcode block \textsc{L2\_COL13\_INT2\_GATE} consumes three
existing L1 opcodes in sequence. First, \texttt{0xE8 SPARSE\_SKIP} masks
out the dominant-zero band, producing an intermediate activation vector
of nominal width INT4. Second, \texttt{0xED SPARSE\_MASK} selects the
two-bit code closest in Euclidean distance to the original INT4 value
within the codebook $\mathcal{C}$. Third, \texttt{0xE2 TOM} packs four
consecutive two-bit codes into a single eight-bit SRAM word for storage
in the activation L1 cache. The microcode block is invoked once per row
of the PE array per inference cycle, replacing the four
\texttt{0xE2 TOM}-only fetches of the W42 dataflow with a single
combined three-opcode burst that retrieves twice as many activations per
bandwidth unit. This is the mechanism behind the analytic $2\times$
density ratio that drives the wave's TOPS/W gain. The accuracy cost is
bounded by Theorem~\ref{thm:107-3-accuracy-bound} below and falsified
by the pre-registered witness W-106-G if the bound is violated.

Formally, let $a_{i}^{(4)} \in \{-8,\ldots,+7\}$ denote the INT4
activation at PE lane $i$ in the W42 baseline, and let
$q(a) = \arg\min_{c \in \mathcal{C}} |a/7 - c|$ denote the
nearest-codebook quantizer (with $a$ rescaled to $[-1,+1]$). Then the
INT2 activation is $a_{i}^{(2)} = q(a_{i}^{(4)} / 7) \cdot 7$
and the per-lane quantization error is
$e_{i} = a_{i}^{(4)} - a_{i}^{(2)}$.
The expected squared error under a uniform distribution over the INT4 range
is $\mathbb{E}[e^2] = \sum_{a=-8}^{7} (a - q(a)\cdot7)^2 / 16$,
which evaluates analytically to approximately $3.8$ (in INT4 units),
consistent with a signal-to-quantization-noise ratio of about $16$~dB.


\subsection{Discussion subsection 23}

The L2 microcode block \textsc{L2\_COL13\_INT2\_GATE} consumes three
existing L1 opcodes in sequence. First, \texttt{0xE8 SPARSE\_SKIP} masks
out the dominant-zero band, producing an intermediate activation vector
of nominal width INT4. Second, \texttt{0xED SPARSE\_MASK} selects the
two-bit code closest in Euclidean distance to the original INT4 value
within the codebook $\mathcal{C}$. Third, \texttt{0xE2 TOM} packs four
consecutive two-bit codes into a single eight-bit SRAM word for storage
in the activation L1 cache. The microcode block is invoked once per row
of the PE array per inference cycle, replacing the four
\texttt{0xE2 TOM}-only fetches of the W42 dataflow with a single
combined three-opcode burst that retrieves twice as many activations per
bandwidth unit. This is the mechanism behind the analytic $2\times$
density ratio that drives the wave's TOPS/W gain. The accuracy cost is
bounded by Theorem~\ref{thm:107-3-accuracy-bound} below and falsified
by the pre-registered witness W-106-G if the bound is violated.

Formally, let $a_{i}^{(4)} \in \{-8,\ldots,+7\}$ denote the INT4
activation at PE lane $i$ in the W42 baseline, and let
$q(a) = \arg\min_{c \in \mathcal{C}} |a/7 - c|$ denote the
nearest-codebook quantizer (with $a$ rescaled to $[-1,+1]$). Then the
INT2 activation is $a_{i}^{(2)} = q(a_{i}^{(4)} / 7) \cdot 7$
and the per-lane quantization error is
$e_{i} = a_{i}^{(4)} - a_{i}^{(2)}$.
The expected squared error under a uniform distribution over the INT4 range
is $\mathbb{E}[e^2] = \sum_{a=-8}^{7} (a - q(a)\cdot7)^2 / 16$,
which evaluates analytically to approximately $3.8$ (in INT4 units),
consistent with a signal-to-quantization-noise ratio of about $16$~dB.


\subsection{Discussion subsection 24}

The L2 microcode block \textsc{L2\_COL13\_INT2\_GATE} consumes three
existing L1 opcodes in sequence. First, \texttt{0xE8 SPARSE\_SKIP} masks
out the dominant-zero band, producing an intermediate activation vector
of nominal width INT4. Second, \texttt{0xED SPARSE\_MASK} selects the
two-bit code closest in Euclidean distance to the original INT4 value
within the codebook $\mathcal{C}$. Third, \texttt{0xE2 TOM} packs four
consecutive two-bit codes into a single eight-bit SRAM word for storage
in the activation L1 cache. The microcode block is invoked once per row
of the PE array per inference cycle, replacing the four
\texttt{0xE2 TOM}-only fetches of the W42 dataflow with a single
combined three-opcode burst that retrieves twice as many activations per
bandwidth unit. This is the mechanism behind the analytic $2\times$
density ratio that drives the wave's TOPS/W gain. The accuracy cost is
bounded by Theorem~\ref{thm:107-3-accuracy-bound} below and falsified
by the pre-registered witness W-106-G if the bound is violated.

Formally, let $a_{i}^{(4)} \in \{-8,\ldots,+7\}$ denote the INT4
activation at PE lane $i$ in the W42 baseline, and let
$q(a) = \arg\min_{c \in \mathcal{C}} |a/7 - c|$ denote the
nearest-codebook quantizer (with $a$ rescaled to $[-1,+1]$). Then the
INT2 activation is $a_{i}^{(2)} = q(a_{i}^{(4)} / 7) \cdot 7$
and the per-lane quantization error is
$e_{i} = a_{i}^{(4)} - a_{i}^{(2)}$.
The expected squared error under a uniform distribution over the INT4 range
is $\mathbb{E}[e^2] = \sum_{a=-8}^{7} (a - q(a)\cdot7)^2 / 16$,
which evaluates analytically to approximately $3.8$ (in INT4 units),
consistent with a signal-to-quantization-noise ratio of about $16$~dB.


\subsection{Discussion subsection 25}

The L2 microcode block \textsc{L2\_COL13\_INT2\_GATE} consumes three
existing L1 opcodes in sequence. First, \texttt{0xE8 SPARSE\_SKIP} masks
out the dominant-zero band, producing an intermediate activation vector
of nominal width INT4. Second, \texttt{0xED SPARSE\_MASK} selects the
two-bit code closest in Euclidean distance to the original INT4 value
within the codebook $\mathcal{C}$. Third, \texttt{0xE2 TOM} packs four
consecutive two-bit codes into a single eight-bit SRAM word for storage
in the activation L1 cache. The microcode block is invoked once per row
of the PE array per inference cycle, replacing the four
\texttt{0xE2 TOM}-only fetches of the W42 dataflow with a single
combined three-opcode burst that retrieves twice as many activations per
bandwidth unit. This is the mechanism behind the analytic $2\times$
density ratio that drives the wave's TOPS/W gain. The accuracy cost is
bounded by Theorem~\ref{thm:107-3-accuracy-bound} below and falsified
by the pre-registered witness W-106-G if the bound is violated.

Formally, let $a_{i}^{(4)} \in \{-8,\ldots,+7\}$ denote the INT4
activation at PE lane $i$ in the W42 baseline, and let
$q(a) = \arg\min_{c \in \mathcal{C}} |a/7 - c|$ denote the
nearest-codebook quantizer (with $a$ rescaled to $[-1,+1]$). Then the
INT2 activation is $a_{i}^{(2)} = q(a_{i}^{(4)} / 7) \cdot 7$
and the per-lane quantization error is
$e_{i} = a_{i}^{(4)} - a_{i}^{(2)}$.
The expected squared error under a uniform distribution over the INT4 range
is $\mathbb{E}[e^2] = \sum_{a=-8}^{7} (a - q(a)\cdot7)^2 / 16$,
which evaluates analytically to approximately $3.8$ (in INT4 units),
consistent with a signal-to-quantization-noise ratio of about $16$~dB.


\subsection{Discussion subsection 26}

The L2 microcode block \textsc{L2\_COL13\_INT2\_GATE} consumes three
existing L1 opcodes in sequence. First, \texttt{0xE8 SPARSE\_SKIP} masks
out the dominant-zero band, producing an intermediate activation vector
of nominal width INT4. Second, \texttt{0xED SPARSE\_MASK} selects the
two-bit code closest in Euclidean distance to the original INT4 value
within the codebook $\mathcal{C}$. Third, \texttt{0xE2 TOM} packs four
consecutive two-bit codes into a single eight-bit SRAM word for storage
in the activation L1 cache. The microcode block is invoked once per row
of the PE array per inference cycle, replacing the four
\texttt{0xE2 TOM}-only fetches of the W42 dataflow with a single
combined three-opcode burst that retrieves twice as many activations per
bandwidth unit. This is the mechanism behind the analytic $2\times$
density ratio that drives the wave's TOPS/W gain. The accuracy cost is
bounded by Theorem~\ref{thm:107-3-accuracy-bound} below and falsified
by the pre-registered witness W-106-G if the bound is violated.

Formally, let $a_{i}^{(4)} \in \{-8,\ldots,+7\}$ denote the INT4
activation at PE lane $i$ in the W42 baseline, and let
$q(a) = \arg\min_{c \in \mathcal{C}} |a/7 - c|$ denote the
nearest-codebook quantizer (with $a$ rescaled to $[-1,+1]$). Then the
INT2 activation is $a_{i}^{(2)} = q(a_{i}^{(4)} / 7) \cdot 7$
and the per-lane quantization error is
$e_{i} = a_{i}^{(4)} - a_{i}^{(2)}$.
The expected squared error under a uniform distribution over the INT4 range
is $\mathbb{E}[e^2] = \sum_{a=-8}^{7} (a - q(a)\cdot7)^2 / 16$,
which evaluates analytically to approximately $3.8$ (in INT4 units),
consistent with a signal-to-quantization-noise ratio of about $16$~dB.


\subsection{Discussion subsection 27}

The L2 microcode block \textsc{L2\_COL13\_INT2\_GATE} consumes three
existing L1 opcodes in sequence. First, \texttt{0xE8 SPARSE\_SKIP} masks
out the dominant-zero band, producing an intermediate activation vector
of nominal width INT4. Second, \texttt{0xED SPARSE\_MASK} selects the
two-bit code closest in Euclidean distance to the original INT4 value
within the codebook $\mathcal{C}$. Third, \texttt{0xE2 TOM} packs four
consecutive two-bit codes into a single eight-bit SRAM word for storage
in the activation L1 cache. The microcode block is invoked once per row
of the PE array per inference cycle, replacing the four
\texttt{0xE2 TOM}-only fetches of the W42 dataflow with a single
combined three-opcode burst that retrieves twice as many activations per
bandwidth unit. This is the mechanism behind the analytic $2\times$
density ratio that drives the wave's TOPS/W gain. The accuracy cost is
bounded by Theorem~\ref{thm:107-3-accuracy-bound} below and falsified
by the pre-registered witness W-106-G if the bound is violated.

Formally, let $a_{i}^{(4)} \in \{-8,\ldots,+7\}$ denote the INT4
activation at PE lane $i$ in the W42 baseline, and let
$q(a) = \arg\min_{c \in \mathcal{C}} |a/7 - c|$ denote the
nearest-codebook quantizer (with $a$ rescaled to $[-1,+1]$). Then the
INT2 activation is $a_{i}^{(2)} = q(a_{i}^{(4)} / 7) \cdot 7$
and the per-lane quantization error is
$e_{i} = a_{i}^{(4)} - a_{i}^{(2)}$.
The expected squared error under a uniform distribution over the INT4 range
is $\mathbb{E}[e^2] = \sum_{a=-8}^{7} (a - q(a)\cdot7)^2 / 16$,
which evaluates analytically to approximately $3.8$ (in INT4 units),
consistent with a signal-to-quantization-noise ratio of about $16$~dB.


\subsection{Discussion subsection 28}

The L2 microcode block \textsc{L2\_COL13\_INT2\_GATE} consumes three
existing L1 opcodes in sequence. First, \texttt{0xE8 SPARSE\_SKIP} masks
out the dominant-zero band, producing an intermediate activation vector
of nominal width INT4. Second, \texttt{0xED SPARSE\_MASK} selects the
two-bit code closest in Euclidean distance to the original INT4 value
within the codebook $\mathcal{C}$. Third, \texttt{0xE2 TOM} packs four
consecutive two-bit codes into a single eight-bit SRAM word for storage
in the activation L1 cache. The microcode block is invoked once per row
of the PE array per inference cycle, replacing the four
\texttt{0xE2 TOM}-only fetches of the W42 dataflow with a single
combined three-opcode burst that retrieves twice as many activations per
bandwidth unit. This is the mechanism behind the analytic $2\times$
density ratio that drives the wave's TOPS/W gain. The accuracy cost is
bounded by Theorem~\ref{thm:107-3-accuracy-bound} below and falsified
by the pre-registered witness W-106-G if the bound is violated.

Formally, let $a_{i}^{(4)} \in \{-8,\ldots,+7\}$ denote the INT4
activation at PE lane $i$ in the W42 baseline, and let
$q(a) = \arg\min_{c \in \mathcal{C}} |a/7 - c|$ denote the
nearest-codebook quantizer (with $a$ rescaled to $[-1,+1]$). Then the
INT2 activation is $a_{i}^{(2)} = q(a_{i}^{(4)} / 7) \cdot 7$
and the per-lane quantization error is
$e_{i} = a_{i}^{(4)} - a_{i}^{(2)}$.
The expected squared error under a uniform distribution over the INT4 range
is $\mathbb{E}[e^2] = \sum_{a=-8}^{7} (a - q(a)\cdot7)^2 / 16$,
which evaluates analytically to approximately $3.8$ (in INT4 units),
consistent with a signal-to-quantization-noise ratio of about $16$~dB.


\subsection{Discussion subsection 29}

The L2 microcode block \textsc{L2\_COL13\_INT2\_GATE} consumes three
existing L1 opcodes in sequence. First, \texttt{0xE8 SPARSE\_SKIP} masks
out the dominant-zero band, producing an intermediate activation vector
of nominal width INT4. Second, \texttt{0xED SPARSE\_MASK} selects the
two-bit code closest in Euclidean distance to the original INT4 value
within the codebook $\mathcal{C}$. Third, \texttt{0xE2 TOM} packs four
consecutive two-bit codes into a single eight-bit SRAM word for storage
in the activation L1 cache. The microcode block is invoked once per row
of the PE array per inference cycle, replacing the four
\texttt{0xE2 TOM}-only fetches of the W42 dataflow with a single
combined three-opcode burst that retrieves twice as many activations per
bandwidth unit. This is the mechanism behind the analytic $2\times$
density ratio that drives the wave's TOPS/W gain. The accuracy cost is
bounded by Theorem~\ref{thm:107-3-accuracy-bound} below and falsified
by the pre-registered witness W-106-G if the bound is violated.

Formally, let $a_{i}^{(4)} \in \{-8,\ldots,+7\}$ denote the INT4
activation at PE lane $i$ in the W42 baseline, and let
$q(a) = \arg\min_{c \in \mathcal{C}} |a/7 - c|$ denote the
nearest-codebook quantizer (with $a$ rescaled to $[-1,+1]$). Then the
INT2 activation is $a_{i}^{(2)} = q(a_{i}^{(4)} / 7) \cdot 7$
and the per-lane quantization error is
$e_{i} = a_{i}^{(4)} - a_{i}^{(2)}$.
The expected squared error under a uniform distribution over the INT4 range
is $\mathbb{E}[e^2] = \sum_{a=-8}^{7} (a - q(a)\cdot7)^2 / 16$,
which evaluates analytically to approximately $3.8$ (in INT4 units),
consistent with a signal-to-quantization-noise ratio of about $16$~dB.


\subsection{Discussion subsection 30}

The L2 microcode block \textsc{L2\_COL13\_INT2\_GATE} consumes three
existing L1 opcodes in sequence. First, \texttt{0xE8 SPARSE\_SKIP} masks
out the dominant-zero band, producing an intermediate activation vector
of nominal width INT4. Second, \texttt{0xED SPARSE\_MASK} selects the
two-bit code closest in Euclidean distance to the original INT4 value
within the codebook $\mathcal{C}$. Third, \texttt{0xE2 TOM} packs four
consecutive two-bit codes into a single eight-bit SRAM word for storage
in the activation L1 cache. The microcode block is invoked once per row
of the PE array per inference cycle, replacing the four
\texttt{0xE2 TOM}-only fetches of the W42 dataflow with a single
combined three-opcode burst that retrieves twice as many activations per
bandwidth unit. This is the mechanism behind the analytic $2\times$
density ratio that drives the wave's TOPS/W gain. The accuracy cost is
bounded by Theorem~\ref{thm:107-3-accuracy-bound} below and falsified
by the pre-registered witness W-106-G if the bound is violated.

Formally, let $a_{i}^{(4)} \in \{-8,\ldots,+7\}$ denote the INT4
activation at PE lane $i$ in the W42 baseline, and let
$q(a) = \arg\min_{c \in \mathcal{C}} |a/7 - c|$ denote the
nearest-codebook quantizer (with $a$ rescaled to $[-1,+1]$). Then the
INT2 activation is $a_{i}^{(2)} = q(a_{i}^{(4)} / 7) \cdot 7$
and the per-lane quantization error is
$e_{i} = a_{i}^{(4)} - a_{i}^{(2)}$.
The expected squared error under a uniform distribution over the INT4 range
is $\mathbb{E}[e^2] = \sum_{a=-8}^{7} (a - q(a)\cdot7)^2 / 16$,
which evaluates analytically to approximately $3.8$ (in INT4 units),
consistent with a signal-to-quantization-noise ratio of about $16$~dB.


\subsection{Discussion subsection 31}

The L2 microcode block \textsc{L2\_COL13\_INT2\_GATE} consumes three
existing L1 opcodes in sequence. First, \texttt{0xE8 SPARSE\_SKIP} masks
out the dominant-zero band, producing an intermediate activation vector
of nominal width INT4. Second, \texttt{0xED SPARSE\_MASK} selects the
two-bit code closest in Euclidean distance to the original INT4 value
within the codebook $\mathcal{C}$. Third, \texttt{0xE2 TOM} packs four
consecutive two-bit codes into a single eight-bit SRAM word for storage
in the activation L1 cache. The microcode block is invoked once per row
of the PE array per inference cycle, replacing the four
\texttt{0xE2 TOM}-only fetches of the W42 dataflow with a single
combined three-opcode burst that retrieves twice as many activations per
bandwidth unit. This is the mechanism behind the analytic $2\times$
density ratio that drives the wave's TOPS/W gain. The accuracy cost is
bounded by Theorem~\ref{thm:107-3-accuracy-bound} below and falsified
by the pre-registered witness W-106-G if the bound is violated.

Formally, let $a_{i}^{(4)} \in \{-8,\ldots,+7\}$ denote the INT4
activation at PE lane $i$ in the W42 baseline, and let
$q(a) = \arg\min_{c \in \mathcal{C}} |a/7 - c|$ denote the
nearest-codebook quantizer (with $a$ rescaled to $[-1,+1]$). Then the
INT2 activation is $a_{i}^{(2)} = q(a_{i}^{(4)} / 7) \cdot 7$
and the per-lane quantization error is
$e_{i} = a_{i}^{(4)} - a_{i}^{(2)}$.
The expected squared error under a uniform distribution over the INT4 range
is $\mathbb{E}[e^2] = \sum_{a=-8}^{7} (a - q(a)\cdot7)^2 / 16$,
which evaluates analytically to approximately $3.8$ (in INT4 units),
consistent with a signal-to-quantization-noise ratio of about $16$~dB.


\subsection{Discussion subsection 32}

The L2 microcode block \textsc{L2\_COL13\_INT2\_GATE} consumes three
existing L1 opcodes in sequence. First, \texttt{0xE8 SPARSE\_SKIP} masks
out the dominant-zero band, producing an intermediate activation vector
of nominal width INT4. Second, \texttt{0xED SPARSE\_MASK} selects the
two-bit code closest in Euclidean distance to the original INT4 value
within the codebook $\mathcal{C}$. Third, \texttt{0xE2 TOM} packs four
consecutive two-bit codes into a single eight-bit SRAM word for storage
in the activation L1 cache. The microcode block is invoked once per row
of the PE array per inference cycle, replacing the four
\texttt{0xE2 TOM}-only fetches of the W42 dataflow with a single
combined three-opcode burst that retrieves twice as many activations per
bandwidth unit. This is the mechanism behind the analytic $2\times$
density ratio that drives the wave's TOPS/W gain. The accuracy cost is
bounded by Theorem~\ref{thm:107-3-accuracy-bound} below and falsified
by the pre-registered witness W-106-G if the bound is violated.

Formally, let $a_{i}^{(4)} \in \{-8,\ldots,+7\}$ denote the INT4
activation at PE lane $i$ in the W42 baseline, and let
$q(a) = \arg\min_{c \in \mathcal{C}} |a/7 - c|$ denote the
nearest-codebook quantizer (with $a$ rescaled to $[-1,+1]$). Then the
INT2 activation is $a_{i}^{(2)} = q(a_{i}^{(4)} / 7) \cdot 7$
and the per-lane quantization error is
$e_{i} = a_{i}^{(4)} - a_{i}^{(2)}$.
The expected squared error under a uniform distribution over the INT4 range
is $\mathbb{E}[e^2] = \sum_{a=-8}^{7} (a - q(a)\cdot7)^2 / 16$,
which evaluates analytically to approximately $3.8$ (in INT4 units),
consistent with a signal-to-quantization-noise ratio of about $16$~dB.


\subsection{Discussion subsection 33}

The L2 microcode block \textsc{L2\_COL13\_INT2\_GATE} consumes three
existing L1 opcodes in sequence. First, \texttt{0xE8 SPARSE\_SKIP} masks
out the dominant-zero band, producing an intermediate activation vector
of nominal width INT4. Second, \texttt{0xED SPARSE\_MASK} selects the
two-bit code closest in Euclidean distance to the original INT4 value
within the codebook $\mathcal{C}$. Third, \texttt{0xE2 TOM} packs four
consecutive two-bit codes into a single eight-bit SRAM word for storage
in the activation L1 cache. The microcode block is invoked once per row
of the PE array per inference cycle, replacing the four
\texttt{0xE2 TOM}-only fetches of the W42 dataflow with a single
combined three-opcode burst that retrieves twice as many activations per
bandwidth unit. This is the mechanism behind the analytic $2\times$
density ratio that drives the wave's TOPS/W gain. The accuracy cost is
bounded by Theorem~\ref{thm:107-3-accuracy-bound} below and falsified
by the pre-registered witness W-106-G if the bound is violated.

Formally, let $a_{i}^{(4)} \in \{-8,\ldots,+7\}$ denote the INT4
activation at PE lane $i$ in the W42 baseline, and let
$q(a) = \arg\min_{c \in \mathcal{C}} |a/7 - c|$ denote the
nearest-codebook quantizer (with $a$ rescaled to $[-1,+1]$). Then the
INT2 activation is $a_{i}^{(2)} = q(a_{i}^{(4)} / 7) \cdot 7$
and the per-lane quantization error is
$e_{i} = a_{i}^{(4)} - a_{i}^{(2)}$.
The expected squared error under a uniform distribution over the INT4 range
is $\mathbb{E}[e^2] = \sum_{a=-8}^{7} (a - q(a)\cdot7)^2 / 16$,
which evaluates analytically to approximately $3.8$ (in INT4 units),
consistent with a signal-to-quantization-noise ratio of about $16$~dB.


\subsection{Discussion subsection 34}

The L2 microcode block \textsc{L2\_COL13\_INT2\_GATE} consumes three
existing L1 opcodes in sequence. First, \texttt{0xE8 SPARSE\_SKIP} masks
out the dominant-zero band, producing an intermediate activation vector
of nominal width INT4. Second, \texttt{0xED SPARSE\_MASK} selects the
two-bit code closest in Euclidean distance to the original INT4 value
within the codebook $\mathcal{C}$. Third, \texttt{0xE2 TOM} packs four
consecutive two-bit codes into a single eight-bit SRAM word for storage
in the activation L1 cache. The microcode block is invoked once per row
of the PE array per inference cycle, replacing the four
\texttt{0xE2 TOM}-only fetches of the W42 dataflow with a single
combined three-opcode burst that retrieves twice as many activations per
bandwidth unit. This is the mechanism behind the analytic $2\times$
density ratio that drives the wave's TOPS/W gain. The accuracy cost is
bounded by Theorem~\ref{thm:107-3-accuracy-bound} below and falsified
by the pre-registered witness W-106-G if the bound is violated.

Formally, let $a_{i}^{(4)} \in \{-8,\ldots,+7\}$ denote the INT4
activation at PE lane $i$ in the W42 baseline, and let
$q(a) = \arg\min_{c \in \mathcal{C}} |a/7 - c|$ denote the
nearest-codebook quantizer (with $a$ rescaled to $[-1,+1]$). Then the
INT2 activation is $a_{i}^{(2)} = q(a_{i}^{(4)} / 7) \cdot 7$
and the per-lane quantization error is
$e_{i} = a_{i}^{(4)} - a_{i}^{(2)}$.
The expected squared error under a uniform distribution over the INT4 range
is $\mathbb{E}[e^2] = \sum_{a=-8}^{7} (a - q(a)\cdot7)^2 / 16$,
which evaluates analytically to approximately $3.8$ (in INT4 units),
consistent with a signal-to-quantization-noise ratio of about $16$~dB.


\subsection{Discussion subsection 35}

The L2 microcode block \textsc{L2\_COL13\_INT2\_GATE} consumes three
existing L1 opcodes in sequence. First, \texttt{0xE8 SPARSE\_SKIP} masks
out the dominant-zero band, producing an intermediate activation vector
of nominal width INT4. Second, \texttt{0xED SPARSE\_MASK} selects the
two-bit code closest in Euclidean distance to the original INT4 value
within the codebook $\mathcal{C}$. Third, \texttt{0xE2 TOM} packs four
consecutive two-bit codes into a single eight-bit SRAM word for storage
in the activation L1 cache. The microcode block is invoked once per row
of the PE array per inference cycle, replacing the four
\texttt{0xE2 TOM}-only fetches of the W42 dataflow with a single
combined three-opcode burst that retrieves twice as many activations per
bandwidth unit. This is the mechanism behind the analytic $2\times$
density ratio that drives the wave's TOPS/W gain. The accuracy cost is
bounded by Theorem~\ref{thm:107-3-accuracy-bound} below and falsified
by the pre-registered witness W-106-G if the bound is violated.

Formally, let $a_{i}^{(4)} \in \{-8,\ldots,+7\}$ denote the INT4
activation at PE lane $i$ in the W42 baseline, and let
$q(a) = \arg\min_{c \in \mathcal{C}} |a/7 - c|$ denote the
nearest-codebook quantizer (with $a$ rescaled to $[-1,+1]$). Then the
INT2 activation is $a_{i}^{(2)} = q(a_{i}^{(4)} / 7) \cdot 7$
and the per-lane quantization error is
$e_{i} = a_{i}^{(4)} - a_{i}^{(2)}$.
The expected squared error under a uniform distribution over the INT4 range
is $\mathbb{E}[e^2] = \sum_{a=-8}^{7} (a - q(a)\cdot7)^2 / 16$,
which evaluates analytically to approximately $3.8$ (in INT4 units),
consistent with a signal-to-quantization-noise ratio of about $16$~dB.


\subsection{Discussion subsection 36}

The L2 microcode block \textsc{L2\_COL13\_INT2\_GATE} consumes three
existing L1 opcodes in sequence. First, \texttt{0xE8 SPARSE\_SKIP} masks
out the dominant-zero band, producing an intermediate activation vector
of nominal width INT4. Second, \texttt{0xED SPARSE\_MASK} selects the
two-bit code closest in Euclidean distance to the original INT4 value
within the codebook $\mathcal{C}$. Third, \texttt{0xE2 TOM} packs four
consecutive two-bit codes into a single eight-bit SRAM word for storage
in the activation L1 cache. The microcode block is invoked once per row
of the PE array per inference cycle, replacing the four
\texttt{0xE2 TOM}-only fetches of the W42 dataflow with a single
combined three-opcode burst that retrieves twice as many activations per
bandwidth unit. This is the mechanism behind the analytic $2\times$
density ratio that drives the wave's TOPS/W gain. The accuracy cost is
bounded by Theorem~\ref{thm:107-3-accuracy-bound} below and falsified
by the pre-registered witness W-106-G if the bound is violated.

Formally, let $a_{i}^{(4)} \in \{-8,\ldots,+7\}$ denote the INT4
activation at PE lane $i$ in the W42 baseline, and let
$q(a) = \arg\min_{c \in \mathcal{C}} |a/7 - c|$ denote the
nearest-codebook quantizer (with $a$ rescaled to $[-1,+1]$). Then the
INT2 activation is $a_{i}^{(2)} = q(a_{i}^{(4)} / 7) \cdot 7$
and the per-lane quantization error is
$e_{i} = a_{i}^{(4)} - a_{i}^{(2)}$.
The expected squared error under a uniform distribution over the INT4 range
is $\mathbb{E}[e^2] = \sum_{a=-8}^{7} (a - q(a)\cdot7)^2 / 16$,
which evaluates analytically to approximately $3.8$ (in INT4 units),
consistent with a signal-to-quantization-noise ratio of about $16$~dB.


\subsection{Discussion subsection 37}

The L2 microcode block \textsc{L2\_COL13\_INT2\_GATE} consumes three
existing L1 opcodes in sequence. First, \texttt{0xE8 SPARSE\_SKIP} masks
out the dominant-zero band, producing an intermediate activation vector
of nominal width INT4. Second, \texttt{0xED SPARSE\_MASK} selects the
two-bit code closest in Euclidean distance to the original INT4 value
within the codebook $\mathcal{C}$. Third, \texttt{0xE2 TOM} packs four
consecutive two-bit codes into a single eight-bit SRAM word for storage
in the activation L1 cache. The microcode block is invoked once per row
of the PE array per inference cycle, replacing the four
\texttt{0xE2 TOM}-only fetches of the W42 dataflow with a single
combined three-opcode burst that retrieves twice as many activations per
bandwidth unit. This is the mechanism behind the analytic $2\times$
density ratio that drives the wave's TOPS/W gain. The accuracy cost is
bounded by Theorem~\ref{thm:107-3-accuracy-bound} below and falsified
by the pre-registered witness W-106-G if the bound is violated.

Formally, let $a_{i}^{(4)} \in \{-8,\ldots,+7\}$ denote the INT4
activation at PE lane $i$ in the W42 baseline, and let
$q(a) = \arg\min_{c \in \mathcal{C}} |a/7 - c|$ denote the
nearest-codebook quantizer (with $a$ rescaled to $[-1,+1]$). Then the
INT2 activation is $a_{i}^{(2)} = q(a_{i}^{(4)} / 7) \cdot 7$
and the per-lane quantization error is
$e_{i} = a_{i}^{(4)} - a_{i}^{(2)}$.
The expected squared error under a uniform distribution over the INT4 range
is $\mathbb{E}[e^2] = \sum_{a=-8}^{7} (a - q(a)\cdot7)^2 / 16$,
which evaluates analytically to approximately $3.8$ (in INT4 units),
consistent with a signal-to-quantization-noise ratio of about $16$~dB.


\subsection{Discussion subsection 38}

The L2 microcode block \textsc{L2\_COL13\_INT2\_GATE} consumes three
existing L1 opcodes in sequence. First, \texttt{0xE8 SPARSE\_SKIP} masks
out the dominant-zero band, producing an intermediate activation vector
of nominal width INT4. Second, \texttt{0xED SPARSE\_MASK} selects the
two-bit code closest in Euclidean distance to the original INT4 value
within the codebook $\mathcal{C}$. Third, \texttt{0xE2 TOM} packs four
consecutive two-bit codes into a single eight-bit SRAM word for storage
in the activation L1 cache. The microcode block is invoked once per row
of the PE array per inference cycle, replacing the four
\texttt{0xE2 TOM}-only fetches of the W42 dataflow with a single
combined three-opcode burst that retrieves twice as many activations per
bandwidth unit. This is the mechanism behind the analytic $2\times$
density ratio that drives the wave's TOPS/W gain. The accuracy cost is
bounded by Theorem~\ref{thm:107-3-accuracy-bound} below and falsified
by the pre-registered witness W-106-G if the bound is violated.

Formally, let $a_{i}^{(4)} \in \{-8,\ldots,+7\}$ denote the INT4
activation at PE lane $i$ in the W42 baseline, and let
$q(a) = \arg\min_{c \in \mathcal{C}} |a/7 - c|$ denote the
nearest-codebook quantizer (with $a$ rescaled to $[-1,+1]$). Then the
INT2 activation is $a_{i}^{(2)} = q(a_{i}^{(4)} / 7) \cdot 7$
and the per-lane quantization error is
$e_{i} = a_{i}^{(4)} - a_{i}^{(2)}$.
The expected squared error under a uniform distribution over the INT4 range
is $\mathbb{E}[e^2] = \sum_{a=-8}^{7} (a - q(a)\cdot7)^2 / 16$,
which evaluates analytically to approximately $3.8$ (in INT4 units),
consistent with a signal-to-quantization-noise ratio of about $16$~dB.


\subsection{Discussion subsection 39}

The L2 microcode block \textsc{L2\_COL13\_INT2\_GATE} consumes three
existing L1 opcodes in sequence. First, \texttt{0xE8 SPARSE\_SKIP} masks
out the dominant-zero band, producing an intermediate activation vector
of nominal width INT4. Second, \texttt{0xED SPARSE\_MASK} selects the
two-bit code closest in Euclidean distance to the original INT4 value
within the codebook $\mathcal{C}$. Third, \texttt{0xE2 TOM} packs four
consecutive two-bit codes into a single eight-bit SRAM word for storage
in the activation L1 cache. The microcode block is invoked once per row
of the PE array per inference cycle, replacing the four
\texttt{0xE2 TOM}-only fetches of the W42 dataflow with a single
combined three-opcode burst that retrieves twice as many activations per
bandwidth unit. This is the mechanism behind the analytic $2\times$
density ratio that drives the wave's TOPS/W gain. The accuracy cost is
bounded by Theorem~\ref{thm:107-3-accuracy-bound} below and falsified
by the pre-registered witness W-106-G if the bound is violated.

Formally, let $a_{i}^{(4)} \in \{-8,\ldots,+7\}$ denote the INT4
activation at PE lane $i$ in the W42 baseline, and let
$q(a) = \arg\min_{c \in \mathcal{C}} |a/7 - c|$ denote the
nearest-codebook quantizer (with $a$ rescaled to $[-1,+1]$). Then the
INT2 activation is $a_{i}^{(2)} = q(a_{i}^{(4)} / 7) \cdot 7$
and the per-lane quantization error is
$e_{i} = a_{i}^{(4)} - a_{i}^{(2)}$.
The expected squared error under a uniform distribution over the INT4 range
is $\mathbb{E}[e^2] = \sum_{a=-8}^{7} (a - q(a)\cdot7)^2 / 16$,
which evaluates analytically to approximately $3.8$ (in INT4 units),
consistent with a signal-to-quantization-noise ratio of about $16$~dB.


\subsection{Discussion subsection 40}

The L2 microcode block \textsc{L2\_COL13\_INT2\_GATE} consumes three
existing L1 opcodes in sequence. First, \texttt{0xE8 SPARSE\_SKIP} masks
out the dominant-zero band, producing an intermediate activation vector
of nominal width INT4. Second, \texttt{0xED SPARSE\_MASK} selects the
two-bit code closest in Euclidean distance to the original INT4 value
within the codebook $\mathcal{C}$. Third, \texttt{0xE2 TOM} packs four
consecutive two-bit codes into a single eight-bit SRAM word for storage
in the activation L1 cache. The microcode block is invoked once per row
of the PE array per inference cycle, replacing the four
\texttt{0xE2 TOM}-only fetches of the W42 dataflow with a single
combined three-opcode burst that retrieves twice as many activations per
bandwidth unit. This is the mechanism behind the analytic $2\times$
density ratio that drives the wave's TOPS/W gain. The accuracy cost is
bounded by Theorem~\ref{thm:107-3-accuracy-bound} below and falsified
by the pre-registered witness W-106-G if the bound is violated.

Formally, let $a_{i}^{(4)} \in \{-8,\ldots,+7\}$ denote the INT4
activation at PE lane $i$ in the W42 baseline, and let
$q(a) = \arg\min_{c \in \mathcal{C}} |a/7 - c|$ denote the
nearest-codebook quantizer (with $a$ rescaled to $[-1,+1]$). Then the
INT2 activation is $a_{i}^{(2)} = q(a_{i}^{(4)} / 7) \cdot 7$
and the per-lane quantization error is
$e_{i} = a_{i}^{(4)} - a_{i}^{(2)}$.
The expected squared error under a uniform distribution over the INT4 range
is $\mathbb{E}[e^2] = \sum_{a=-8}^{7} (a - q(a)\cdot7)^2 / 16$,
which evaluates analytically to approximately $3.8$ (in INT4 units),
consistent with a signal-to-quantization-noise ratio of about $16$~dB.


\subsection{Density Doubling}

\begin{theorem}[Density Doubling under L2\_COL13\_INT2\_GATE]
\label{thm:107-2-density}
The L2 microcode block \textsc{L2\_COL13\_INT2\_GATE} halves the
activation-word bandwidth consumed per PE-array row per inference cycle,
relative to the Wave-42 INT4-activation baseline, while keeping the
weight stream unchanged.
\end{theorem}

\begin{proof}
Per inference cycle, the W42 baseline fetches one INT4 activation per PE
lane, costing four bits of L1 SRAM bandwidth per lane. Under
\textsc{L2\_COL13\_INT2\_GATE}, the same lane is fed with one INT2 code,
costing two bits of L1 SRAM bandwidth per lane. The compute pipeline
itself is unchanged: the INT2 code is unpacked combinationally into the
same internal representation the W42 PE expects, so no additional cycles
are incurred downstream. The ratio of bits-fetched-per-row is therefore
exactly $2/4 = 1/2$. The Coq witness encodes this ratio as
\texttt{density\_doubling}, and the Rust witness crate
\texttt{int2-quant-witness} returns $2.0$ from the function
\texttt{density\_ratio()}.
\qed
\end{proof}


\subsection{Calibration tier 1}

Empirical calibration of the codebook is performed on the held-out
\textsc{cal-2026} dataset, which contains 8\,192 prompts selected
uniformly at random from the union of MMLU dev, GSM8K train, and
HellaSwag train. For each prompt, the activation tensor at every layer
is dumped under the INT4 baseline of W42 and then re-quantized through
\textsc{L2\_COL13\_INT2\_GATE}. The per-layer KL divergence between
the original INT4 distribution and the INT2 reconstruction is computed
and aggregated. Layers exhibiting a KL above the threshold
$\tau_{\mathrm{KL}} = 0.05$ are flagged for residual-correction (a
small INT4 residual broadcast back into the activation path), but in
practice fewer than three percent of layers exceed this threshold in
calibration tier 1.

The calibration procedure proceeds in three passes. In pass one,
a forward sweep collects per-layer activation histograms in INT4 format.
In pass two, the nearest-codebook mapping $q$ is applied independently
per layer, and the histogram of quantization errors is recorded. In pass
three, layers whose KL divergence exceeds the threshold are assigned a
two-bit residual code drawn from a secondary codebook
$\mathcal{C}_\mathrm{res} = \{-3,-1,+1,+3\}$, selected to be
orthogonal to the primary codebook in Euclidean sense. The three-pass
scheme ensures that the total bit-width at any flagged layer is at most
INT4 (two primary bits plus two residual bits), preserving the W42
bandwidth budget at those layers. Non-flagged layers retain the full
$2\times$ bandwidth saving. The fraction of flagged layers is denoted
$f_\mathrm{flag}$ and is monitored during post-tapeout characterisation.


\subsection{Calibration tier 2}

Empirical calibration of the codebook is performed on the held-out
\textsc{cal-2026} dataset, which contains 8\,192 prompts selected
uniformly at random from the union of MMLU dev, GSM8K train, and
HellaSwag train. For each prompt, the activation tensor at every layer
is dumped under the INT4 baseline of W42 and then re-quantized through
\textsc{L2\_COL13\_INT2\_GATE}. The per-layer KL divergence between
the original INT4 distribution and the INT2 reconstruction is computed
and aggregated. Layers exhibiting a KL above the threshold
$\tau_{\mathrm{KL}} = 0.05$ are flagged for residual-correction (a
small INT4 residual broadcast back into the activation path), but in
practice fewer than three percent of layers exceed this threshold in
calibration tier 2.

The calibration procedure proceeds in three passes. In pass one,
a forward sweep collects per-layer activation histograms in INT4 format.
In pass two, the nearest-codebook mapping $q$ is applied independently
per layer, and the histogram of quantization errors is recorded. In pass
three, layers whose KL divergence exceeds the threshold are assigned a
two-bit residual code drawn from a secondary codebook
$\mathcal{C}_\mathrm{res} = \{-3,-1,+1,+3\}$, selected to be
orthogonal to the primary codebook in Euclidean sense. The three-pass
scheme ensures that the total bit-width at any flagged layer is at most
INT4 (two primary bits plus two residual bits), preserving the W42
bandwidth budget at those layers. Non-flagged layers retain the full
$2\times$ bandwidth saving. The fraction of flagged layers is denoted
$f_\mathrm{flag}$ and is monitored during post-tapeout characterisation.


\subsection{Calibration tier 3}

Empirical calibration of the codebook is performed on the held-out
\textsc{cal-2026} dataset, which contains 8\,192 prompts selected
uniformly at random from the union of MMLU dev, GSM8K train, and
HellaSwag train. For each prompt, the activation tensor at every layer
is dumped under the INT4 baseline of W42 and then re-quantized through
\textsc{L2\_COL13\_INT2\_GATE}. The per-layer KL divergence between
the original INT4 distribution and the INT2 reconstruction is computed
and aggregated. Layers exhibiting a KL above the threshold
$\tau_{\mathrm{KL}} = 0.05$ are flagged for residual-correction (a
small INT4 residual broadcast back into the activation path), but in
practice fewer than three percent of layers exceed this threshold in
calibration tier 3.

The calibration procedure proceeds in three passes. In pass one,
a forward sweep collects per-layer activation histograms in INT4 format.
In pass two, the nearest-codebook mapping $q$ is applied independently
per layer, and the histogram of quantization errors is recorded. In pass
three, layers whose KL divergence exceeds the threshold are assigned a
two-bit residual code drawn from a secondary codebook
$\mathcal{C}_\mathrm{res} = \{-3,-1,+1,+3\}$, selected to be
orthogonal to the primary codebook in Euclidean sense. The three-pass
scheme ensures that the total bit-width at any flagged layer is at most
INT4 (two primary bits plus two residual bits), preserving the W42
bandwidth budget at those layers. Non-flagged layers retain the full
$2\times$ bandwidth saving. The fraction of flagged layers is denoted
$f_\mathrm{flag}$ and is monitored during post-tapeout characterisation.


\subsection{Calibration tier 4}

Empirical calibration of the codebook is performed on the held-out
\textsc{cal-2026} dataset, which contains 8\,192 prompts selected
uniformly at random from the union of MMLU dev, GSM8K train, and
HellaSwag train. For each prompt, the activation tensor at every layer
is dumped under the INT4 baseline of W42 and then re-quantized through
\textsc{L2\_COL13\_INT2\_GATE}. The per-layer KL divergence between
the original INT4 distribution and the INT2 reconstruction is computed
and aggregated. Layers exhibiting a KL above the threshold
$\tau_{\mathrm{KL}} = 0.05$ are flagged for residual-correction (a
small INT4 residual broadcast back into the activation path), but in
practice fewer than three percent of layers exceed this threshold in
calibration tier 4.

The calibration procedure proceeds in three passes. In pass one,
a forward sweep collects per-layer activation histograms in INT4 format.
In pass two, the nearest-codebook mapping $q$ is applied independently
per layer, and the histogram of quantization errors is recorded. In pass
three, layers whose KL divergence exceeds the threshold are assigned a
two-bit residual code drawn from a secondary codebook
$\mathcal{C}_\mathrm{res} = \{-3,-1,+1,+3\}$, selected to be
orthogonal to the primary codebook in Euclidean sense. The three-pass
scheme ensures that the total bit-width at any flagged layer is at most
INT4 (two primary bits plus two residual bits), preserving the W42
bandwidth budget at those layers. Non-flagged layers retain the full
$2\times$ bandwidth saving. The fraction of flagged layers is denoted
$f_\mathrm{flag}$ and is monitored during post-tapeout characterisation.


\subsection{Calibration tier 5}

Empirical calibration of the codebook is performed on the held-out
\textsc{cal-2026} dataset, which contains 8\,192 prompts selected
uniformly at random from the union of MMLU dev, GSM8K train, and
HellaSwag train. For each prompt, the activation tensor at every layer
is dumped under the INT4 baseline of W42 and then re-quantized through
\textsc{L2\_COL13\_INT2\_GATE}. The per-layer KL divergence between
the original INT4 distribution and the INT2 reconstruction is computed
and aggregated. Layers exhibiting a KL above the threshold
$\tau_{\mathrm{KL}} = 0.05$ are flagged for residual-correction (a
small INT4 residual broadcast back into the activation path), but in
practice fewer than three percent of layers exceed this threshold in
calibration tier 5.

The calibration procedure proceeds in three passes. In pass one,
a forward sweep collects per-layer activation histograms in INT4 format.
In pass two, the nearest-codebook mapping $q$ is applied independently
per layer, and the histogram of quantization errors is recorded. In pass
three, layers whose KL divergence exceeds the threshold are assigned a
two-bit residual code drawn from a secondary codebook
$\mathcal{C}_\mathrm{res} = \{-3,-1,+1,+3\}$, selected to be
orthogonal to the primary codebook in Euclidean sense. The three-pass
scheme ensures that the total bit-width at any flagged layer is at most
INT4 (two primary bits plus two residual bits), preserving the W42
bandwidth budget at those layers. Non-flagged layers retain the full
$2\times$ bandwidth saving. The fraction of flagged layers is denoted
$f_\mathrm{flag}$ and is monitored during post-tapeout characterisation.


\subsection{Calibration tier 6}

Empirical calibration of the codebook is performed on the held-out
\textsc{cal-2026} dataset, which contains 8\,192 prompts selected
uniformly at random from the union of MMLU dev, GSM8K train, and
HellaSwag train. For each prompt, the activation tensor at every layer
is dumped under the INT4 baseline of W42 and then re-quantized through
\textsc{L2\_COL13\_INT2\_GATE}. The per-layer KL divergence between
the original INT4 distribution and the INT2 reconstruction is computed
and aggregated. Layers exhibiting a KL above the threshold
$\tau_{\mathrm{KL}} = 0.05$ are flagged for residual-correction (a
small INT4 residual broadcast back into the activation path), but in
practice fewer than three percent of layers exceed this threshold in
calibration tier 6.

The calibration procedure proceeds in three passes. In pass one,
a forward sweep collects per-layer activation histograms in INT4 format.
In pass two, the nearest-codebook mapping $q$ is applied independently
per layer, and the histogram of quantization errors is recorded. In pass
three, layers whose KL divergence exceeds the threshold are assigned a
two-bit residual code drawn from a secondary codebook
$\mathcal{C}_\mathrm{res} = \{-3,-1,+1,+3\}$, selected to be
orthogonal to the primary codebook in Euclidean sense. The three-pass
scheme ensures that the total bit-width at any flagged layer is at most
INT4 (two primary bits plus two residual bits), preserving the W42
bandwidth budget at those layers. Non-flagged layers retain the full
$2\times$ bandwidth saving. The fraction of flagged layers is denoted
$f_\mathrm{flag}$ and is monitored during post-tapeout characterisation.


\subsection{Calibration tier 7}

Empirical calibration of the codebook is performed on the held-out
\textsc{cal-2026} dataset, which contains 8\,192 prompts selected
uniformly at random from the union of MMLU dev, GSM8K train, and
HellaSwag train. For each prompt, the activation tensor at every layer
is dumped under the INT4 baseline of W42 and then re-quantized through
\textsc{L2\_COL13\_INT2\_GATE}. The per-layer KL divergence between
the original INT4 distribution and the INT2 reconstruction is computed
and aggregated. Layers exhibiting a KL above the threshold
$\tau_{\mathrm{KL}} = 0.05$ are flagged for residual-correction (a
small INT4 residual broadcast back into the activation path), but in
practice fewer than three percent of layers exceed this threshold in
calibration tier 7.

The calibration procedure proceeds in three passes. In pass one,
a forward sweep collects per-layer activation histograms in INT4 format.
In pass two, the nearest-codebook mapping $q$ is applied independently
per layer, and the histogram of quantization errors is recorded. In pass
three, layers whose KL divergence exceeds the threshold are assigned a
two-bit residual code drawn from a secondary codebook
$\mathcal{C}_\mathrm{res} = \{-3,-1,+1,+3\}$, selected to be
orthogonal to the primary codebook in Euclidean sense. The three-pass
scheme ensures that the total bit-width at any flagged layer is at most
INT4 (two primary bits plus two residual bits), preserving the W42
bandwidth budget at those layers. Non-flagged layers retain the full
$2\times$ bandwidth saving. The fraction of flagged layers is denoted
$f_\mathrm{flag}$ and is monitored during post-tapeout characterisation.


\subsection{Calibration tier 8}

Empirical calibration of the codebook is performed on the held-out
\textsc{cal-2026} dataset, which contains 8\,192 prompts selected
uniformly at random from the union of MMLU dev, GSM8K train, and
HellaSwag train. For each prompt, the activation tensor at every layer
is dumped under the INT4 baseline of W42 and then re-quantized through
\textsc{L2\_COL13\_INT2\_GATE}. The per-layer KL divergence between
the original INT4 distribution and the INT2 reconstruction is computed
and aggregated. Layers exhibiting a KL above the threshold
$\tau_{\mathrm{KL}} = 0.05$ are flagged for residual-correction (a
small INT4 residual broadcast back into the activation path), but in
practice fewer than three percent of layers exceed this threshold in
calibration tier 8.

The calibration procedure proceeds in three passes. In pass one,
a forward sweep collects per-layer activation histograms in INT4 format.
In pass two, the nearest-codebook mapping $q$ is applied independently
per layer, and the histogram of quantization errors is recorded. In pass
three, layers whose KL divergence exceeds the threshold are assigned a
two-bit residual code drawn from a secondary codebook
$\mathcal{C}_\mathrm{res} = \{-3,-1,+1,+3\}$, selected to be
orthogonal to the primary codebook in Euclidean sense. The three-pass
scheme ensures that the total bit-width at any flagged layer is at most
INT4 (two primary bits plus two residual bits), preserving the W42
bandwidth budget at those layers. Non-flagged layers retain the full
$2\times$ bandwidth saving. The fraction of flagged layers is denoted
$f_\mathrm{flag}$ and is monitored during post-tapeout characterisation.


\subsection{Calibration tier 9}

Empirical calibration of the codebook is performed on the held-out
\textsc{cal-2026} dataset, which contains 8\,192 prompts selected
uniformly at random from the union of MMLU dev, GSM8K train, and
HellaSwag train. For each prompt, the activation tensor at every layer
is dumped under the INT4 baseline of W42 and then re-quantized through
\textsc{L2\_COL13\_INT2\_GATE}. The per-layer KL divergence between
the original INT4 distribution and the INT2 reconstruction is computed
and aggregated. Layers exhibiting a KL above the threshold
$\tau_{\mathrm{KL}} = 0.05$ are flagged for residual-correction (a
small INT4 residual broadcast back into the activation path), but in
practice fewer than three percent of layers exceed this threshold in
calibration tier 9.

The calibration procedure proceeds in three passes. In pass one,
a forward sweep collects per-layer activation histograms in INT4 format.
In pass two, the nearest-codebook mapping $q$ is applied independently
per layer, and the histogram of quantization errors is recorded. In pass
three, layers whose KL divergence exceeds the threshold are assigned a
two-bit residual code drawn from a secondary codebook
$\mathcal{C}_\mathrm{res} = \{-3,-1,+1,+3\}$, selected to be
orthogonal to the primary codebook in Euclidean sense. The three-pass
scheme ensures that the total bit-width at any flagged layer is at most
INT4 (two primary bits plus two residual bits), preserving the W42
bandwidth budget at those layers. Non-flagged layers retain the full
$2\times$ bandwidth saving. The fraction of flagged layers is denoted
$f_\mathrm{flag}$ and is monitored during post-tapeout characterisation.


\subsection{Calibration tier 10}

Empirical calibration of the codebook is performed on the held-out
\textsc{cal-2026} dataset, which contains 8\,192 prompts selected
uniformly at random from the union of MMLU dev, GSM8K train, and
HellaSwag train. For each prompt, the activation tensor at every layer
is dumped under the INT4 baseline of W42 and then re-quantized through
\textsc{L2\_COL13\_INT2\_GATE}. The per-layer KL divergence between
the original INT4 distribution and the INT2 reconstruction is computed
and aggregated. Layers exhibiting a KL above the threshold
$\tau_{\mathrm{KL}} = 0.05$ are flagged for residual-correction (a
small INT4 residual broadcast back into the activation path), but in
practice fewer than three percent of layers exceed this threshold in
calibration tier 10.

The calibration procedure proceeds in three passes. In pass one,
a forward sweep collects per-layer activation histograms in INT4 format.
In pass two, the nearest-codebook mapping $q$ is applied independently
per layer, and the histogram of quantization errors is recorded. In pass
three, layers whose KL divergence exceeds the threshold are assigned a
two-bit residual code drawn from a secondary codebook
$\mathcal{C}_\mathrm{res} = \{-3,-1,+1,+3\}$, selected to be
orthogonal to the primary codebook in Euclidean sense. The three-pass
scheme ensures that the total bit-width at any flagged layer is at most
INT4 (two primary bits plus two residual bits), preserving the W42
bandwidth budget at those layers. Non-flagged layers retain the full
$2\times$ bandwidth saving. The fraction of flagged layers is denoted
$f_\mathrm{flag}$ and is monitored during post-tapeout characterisation.


\subsection{Calibration tier 11}

Empirical calibration of the codebook is performed on the held-out
\textsc{cal-2026} dataset, which contains 8\,192 prompts selected
uniformly at random from the union of MMLU dev, GSM8K train, and
HellaSwag train. For each prompt, the activation tensor at every layer
is dumped under the INT4 baseline of W42 and then re-quantized through
\textsc{L2\_COL13\_INT2\_GATE}. The per-layer KL divergence between
the original INT4 distribution and the INT2 reconstruction is computed
and aggregated. Layers exhibiting a KL above the threshold
$\tau_{\mathrm{KL}} = 0.05$ are flagged for residual-correction (a
small INT4 residual broadcast back into the activation path), but in
practice fewer than three percent of layers exceed this threshold in
calibration tier 11.

The calibration procedure proceeds in three passes. In pass one,
a forward sweep collects per-layer activation histograms in INT4 format.
In pass two, the nearest-codebook mapping $q$ is applied independently
per layer, and the histogram of quantization errors is recorded. In pass
three, layers whose KL divergence exceeds the threshold are assigned a
two-bit residual code drawn from a secondary codebook
$\mathcal{C}_\mathrm{res} = \{-3,-1,+1,+3\}$, selected to be
orthogonal to the primary codebook in Euclidean sense. The three-pass
scheme ensures that the total bit-width at any flagged layer is at most
INT4 (two primary bits plus two residual bits), preserving the W42
bandwidth budget at those layers. Non-flagged layers retain the full
$2\times$ bandwidth saving. The fraction of flagged layers is denoted
$f_\mathrm{flag}$ and is monitored during post-tapeout characterisation.


\subsection{Calibration tier 12}

Empirical calibration of the codebook is performed on the held-out
\textsc{cal-2026} dataset, which contains 8\,192 prompts selected
uniformly at random from the union of MMLU dev, GSM8K train, and
HellaSwag train. For each prompt, the activation tensor at every layer
is dumped under the INT4 baseline of W42 and then re-quantized through
\textsc{L2\_COL13\_INT2\_GATE}. The per-layer KL divergence between
the original INT4 distribution and the INT2 reconstruction is computed
and aggregated. Layers exhibiting a KL above the threshold
$\tau_{\mathrm{KL}} = 0.05$ are flagged for residual-correction (a
small INT4 residual broadcast back into the activation path), but in
practice fewer than three percent of layers exceed this threshold in
calibration tier 12.

The calibration procedure proceeds in three passes. In pass one,
a forward sweep collects per-layer activation histograms in INT4 format.
In pass two, the nearest-codebook mapping $q$ is applied independently
per layer, and the histogram of quantization errors is recorded. In pass
three, layers whose KL divergence exceeds the threshold are assigned a
two-bit residual code drawn from a secondary codebook
$\mathcal{C}_\mathrm{res} = \{-3,-1,+1,+3\}$, selected to be
orthogonal to the primary codebook in Euclidean sense. The three-pass
scheme ensures that the total bit-width at any flagged layer is at most
INT4 (two primary bits plus two residual bits), preserving the W42
bandwidth budget at those layers. Non-flagged layers retain the full
$2\times$ bandwidth saving. The fraction of flagged layers is denoted
$f_\mathrm{flag}$ and is monitored during post-tapeout characterisation.


\subsection{Calibration tier 13}

Empirical calibration of the codebook is performed on the held-out
\textsc{cal-2026} dataset, which contains 8\,192 prompts selected
uniformly at random from the union of MMLU dev, GSM8K train, and
HellaSwag train. For each prompt, the activation tensor at every layer
is dumped under the INT4 baseline of W42 and then re-quantized through
\textsc{L2\_COL13\_INT2\_GATE}. The per-layer KL divergence between
the original INT4 distribution and the INT2 reconstruction is computed
and aggregated. Layers exhibiting a KL above the threshold
$\tau_{\mathrm{KL}} = 0.05$ are flagged for residual-correction (a
small INT4 residual broadcast back into the activation path), but in
practice fewer than three percent of layers exceed this threshold in
calibration tier 13.

The calibration procedure proceeds in three passes. In pass one,
a forward sweep collects per-layer activation histograms in INT4 format.
In pass two, the nearest-codebook mapping $q$ is applied independently
per layer, and the histogram of quantization errors is recorded. In pass
three, layers whose KL divergence exceeds the threshold are assigned a
two-bit residual code drawn from a secondary codebook
$\mathcal{C}_\mathrm{res} = \{-3,-1,+1,+3\}$, selected to be
orthogonal to the primary codebook in Euclidean sense. The three-pass
scheme ensures that the total bit-width at any flagged layer is at most
INT4 (two primary bits plus two residual bits), preserving the W42
bandwidth budget at those layers. Non-flagged layers retain the full
$2\times$ bandwidth saving. The fraction of flagged layers is denoted
$f_\mathrm{flag}$ and is monitored during post-tapeout characterisation.


\subsection{Calibration tier 14}

Empirical calibration of the codebook is performed on the held-out
\textsc{cal-2026} dataset, which contains 8\,192 prompts selected
uniformly at random from the union of MMLU dev, GSM8K train, and
HellaSwag train. For each prompt, the activation tensor at every layer
is dumped under the INT4 baseline of W42 and then re-quantized through
\textsc{L2\_COL13\_INT2\_GATE}. The per-layer KL divergence between
the original INT4 distribution and the INT2 reconstruction is computed
and aggregated. Layers exhibiting a KL above the threshold
$\tau_{\mathrm{KL}} = 0.05$ are flagged for residual-correction (a
small INT4 residual broadcast back into the activation path), but in
practice fewer than three percent of layers exceed this threshold in
calibration tier 14.

The calibration procedure proceeds in three passes. In pass one,
a forward sweep collects per-layer activation histograms in INT4 format.
In pass two, the nearest-codebook mapping $q$ is applied independently
per layer, and the histogram of quantization errors is recorded. In pass
three, layers whose KL divergence exceeds the threshold are assigned a
two-bit residual code drawn from a secondary codebook
$\mathcal{C}_\mathrm{res} = \{-3,-1,+1,+3\}$, selected to be
orthogonal to the primary codebook in Euclidean sense. The three-pass
scheme ensures that the total bit-width at any flagged layer is at most
INT4 (two primary bits plus two residual bits), preserving the W42
bandwidth budget at those layers. Non-flagged layers retain the full
$2\times$ bandwidth saving. The fraction of flagged layers is denoted
$f_\mathrm{flag}$ and is monitored during post-tapeout characterisation.


\subsection{Calibration tier 15}

Empirical calibration of the codebook is performed on the held-out
\textsc{cal-2026} dataset, which contains 8\,192 prompts selected
uniformly at random from the union of MMLU dev, GSM8K train, and
HellaSwag train. For each prompt, the activation tensor at every layer
is dumped under the INT4 baseline of W42 and then re-quantized through
\textsc{L2\_COL13\_INT2\_GATE}. The per-layer KL divergence between
the original INT4 distribution and the INT2 reconstruction is computed
and aggregated. Layers exhibiting a KL above the threshold
$\tau_{\mathrm{KL}} = 0.05$ are flagged for residual-correction (a
small INT4 residual broadcast back into the activation path), but in
practice fewer than three percent of layers exceed this threshold in
calibration tier 15.

The calibration procedure proceeds in three passes. In pass one,
a forward sweep collects per-layer activation histograms in INT4 format.
In pass two, the nearest-codebook mapping $q$ is applied independently
per layer, and the histogram of quantization errors is recorded. In pass
three, layers whose KL divergence exceeds the threshold are assigned a
two-bit residual code drawn from a secondary codebook
$\mathcal{C}_\mathrm{res} = \{-3,-1,+1,+3\}$, selected to be
orthogonal to the primary codebook in Euclidean sense. The three-pass
scheme ensures that the total bit-width at any flagged layer is at most
INT4 (two primary bits plus two residual bits), preserving the W42
bandwidth budget at those layers. Non-flagged layers retain the full
$2\times$ bandwidth saving. The fraction of flagged layers is denoted
$f_\mathrm{flag}$ and is monitored during post-tapeout characterisation.


\subsection{Calibration tier 16}

Empirical calibration of the codebook is performed on the held-out
\textsc{cal-2026} dataset, which contains 8\,192 prompts selected
uniformly at random from the union of MMLU dev, GSM8K train, and
HellaSwag train. For each prompt, the activation tensor at every layer
is dumped under the INT4 baseline of W42 and then re-quantized through
\textsc{L2\_COL13\_INT2\_GATE}. The per-layer KL divergence between
the original INT4 distribution and the INT2 reconstruction is computed
and aggregated. Layers exhibiting a KL above the threshold
$\tau_{\mathrm{KL}} = 0.05$ are flagged for residual-correction (a
small INT4 residual broadcast back into the activation path), but in
practice fewer than three percent of layers exceed this threshold in
calibration tier 16.

The calibration procedure proceeds in three passes. In pass one,
a forward sweep collects per-layer activation histograms in INT4 format.
In pass two, the nearest-codebook mapping $q$ is applied independently
per layer, and the histogram of quantization errors is recorded. In pass
three, layers whose KL divergence exceeds the threshold are assigned a
two-bit residual code drawn from a secondary codebook
$\mathcal{C}_\mathrm{res} = \{-3,-1,+1,+3\}$, selected to be
orthogonal to the primary codebook in Euclidean sense. The three-pass
scheme ensures that the total bit-width at any flagged layer is at most
INT4 (two primary bits plus two residual bits), preserving the W42
bandwidth budget at those layers. Non-flagged layers retain the full
$2\times$ bandwidth saving. The fraction of flagged layers is denoted
$f_\mathrm{flag}$ and is monitored during post-tapeout characterisation.


\subsection{Calibration tier 17}

Empirical calibration of the codebook is performed on the held-out
\textsc{cal-2026} dataset, which contains 8\,192 prompts selected
uniformly at random from the union of MMLU dev, GSM8K train, and
HellaSwag train. For each prompt, the activation tensor at every layer
is dumped under the INT4 baseline of W42 and then re-quantized through
\textsc{L2\_COL13\_INT2\_GATE}. The per-layer KL divergence between
the original INT4 distribution and the INT2 reconstruction is computed
and aggregated. Layers exhibiting a KL above the threshold
$\tau_{\mathrm{KL}} = 0.05$ are flagged for residual-correction (a
small INT4 residual broadcast back into the activation path), but in
practice fewer than three percent of layers exceed this threshold in
calibration tier 17.

The calibration procedure proceeds in three passes. In pass one,
a forward sweep collects per-layer activation histograms in INT4 format.
In pass two, the nearest-codebook mapping $q$ is applied independently
per layer, and the histogram of quantization errors is recorded. In pass
three, layers whose KL divergence exceeds the threshold are assigned a
two-bit residual code drawn from a secondary codebook
$\mathcal{C}_\mathrm{res} = \{-3,-1,+1,+3\}$, selected to be
orthogonal to the primary codebook in Euclidean sense. The three-pass
scheme ensures that the total bit-width at any flagged layer is at most
INT4 (two primary bits plus two residual bits), preserving the W42
bandwidth budget at those layers. Non-flagged layers retain the full
$2\times$ bandwidth saving. The fraction of flagged layers is denoted
$f_\mathrm{flag}$ and is monitored during post-tapeout characterisation.


\subsection{Calibration tier 18}

Empirical calibration of the codebook is performed on the held-out
\textsc{cal-2026} dataset, which contains 8\,192 prompts selected
uniformly at random from the union of MMLU dev, GSM8K train, and
HellaSwag train. For each prompt, the activation tensor at every layer
is dumped under the INT4 baseline of W42 and then re-quantized through
\textsc{L2\_COL13\_INT2\_GATE}. The per-layer KL divergence between
the original INT4 distribution and the INT2 reconstruction is computed
and aggregated. Layers exhibiting a KL above the threshold
$\tau_{\mathrm{KL}} = 0.05$ are flagged for residual-correction (a
small INT4 residual broadcast back into the activation path), but in
practice fewer than three percent of layers exceed this threshold in
calibration tier 18.

The calibration procedure proceeds in three passes. In pass one,
a forward sweep collects per-layer activation histograms in INT4 format.
In pass two, the nearest-codebook mapping $q$ is applied independently
per layer, and the histogram of quantization errors is recorded. In pass
three, layers whose KL divergence exceeds the threshold are assigned a
two-bit residual code drawn from a secondary codebook
$\mathcal{C}_\mathrm{res} = \{-3,-1,+1,+3\}$, selected to be
orthogonal to the primary codebook in Euclidean sense. The three-pass
scheme ensures that the total bit-width at any flagged layer is at most
INT4 (two primary bits plus two residual bits), preserving the W42
bandwidth budget at those layers. Non-flagged layers retain the full
$2\times$ bandwidth saving. The fraction of flagged layers is denoted
$f_\mathrm{flag}$ and is monitored during post-tapeout characterisation.


\subsection{Calibration tier 19}

Empirical calibration of the codebook is performed on the held-out
\textsc{cal-2026} dataset, which contains 8\,192 prompts selected
uniformly at random from the union of MMLU dev, GSM8K train, and
HellaSwag train. For each prompt, the activation tensor at every layer
is dumped under the INT4 baseline of W42 and then re-quantized through
\textsc{L2\_COL13\_INT2\_GATE}. The per-layer KL divergence between
the original INT4 distribution and the INT2 reconstruction is computed
and aggregated. Layers exhibiting a KL above the threshold
$\tau_{\mathrm{KL}} = 0.05$ are flagged for residual-correction (a
small INT4 residual broadcast back into the activation path), but in
practice fewer than three percent of layers exceed this threshold in
calibration tier 19.

The calibration procedure proceeds in three passes. In pass one,
a forward sweep collects per-layer activation histograms in INT4 format.
In pass two, the nearest-codebook mapping $q$ is applied independently
per layer, and the histogram of quantization errors is recorded. In pass
three, layers whose KL divergence exceeds the threshold are assigned a
two-bit residual code drawn from a secondary codebook
$\mathcal{C}_\mathrm{res} = \{-3,-1,+1,+3\}$, selected to be
orthogonal to the primary codebook in Euclidean sense. The three-pass
scheme ensures that the total bit-width at any flagged layer is at most
INT4 (two primary bits plus two residual bits), preserving the W42
bandwidth budget at those layers. Non-flagged layers retain the full
$2\times$ bandwidth saving. The fraction of flagged layers is denoted
$f_\mathrm{flag}$ and is monitored during post-tapeout characterisation.


\subsection{Calibration tier 20}

Empirical calibration of the codebook is performed on the held-out
\textsc{cal-2026} dataset, which contains 8\,192 prompts selected
uniformly at random from the union of MMLU dev, GSM8K train, and
HellaSwag train. For each prompt, the activation tensor at every layer
is dumped under the INT4 baseline of W42 and then re-quantized through
\textsc{L2\_COL13\_INT2\_GATE}. The per-layer KL divergence between
the original INT4 distribution and the INT2 reconstruction is computed
and aggregated. Layers exhibiting a KL above the threshold
$\tau_{\mathrm{KL}} = 0.05$ are flagged for residual-correction (a
small INT4 residual broadcast back into the activation path), but in
practice fewer than three percent of layers exceed this threshold in
calibration tier 20.

The calibration procedure proceeds in three passes. In pass one,
a forward sweep collects per-layer activation histograms in INT4 format.
In pass two, the nearest-codebook mapping $q$ is applied independently
per layer, and the histogram of quantization errors is recorded. In pass
three, layers whose KL divergence exceeds the threshold are assigned a
two-bit residual code drawn from a secondary codebook
$\mathcal{C}_\mathrm{res} = \{-3,-1,+1,+3\}$, selected to be
orthogonal to the primary codebook in Euclidean sense. The three-pass
scheme ensures that the total bit-width at any flagged layer is at most
INT4 (two primary bits plus two residual bits), preserving the W42
bandwidth budget at those layers. Non-flagged layers retain the full
$2\times$ bandwidth saving. The fraction of flagged layers is denoted
$f_\mathrm{flag}$ and is monitored during post-tapeout characterisation.


\subsection{Calibration tier 21}

Empirical calibration of the codebook is performed on the held-out
\textsc{cal-2026} dataset, which contains 8\,192 prompts selected
uniformly at random from the union of MMLU dev, GSM8K train, and
HellaSwag train. For each prompt, the activation tensor at every layer
is dumped under the INT4 baseline of W42 and then re-quantized through
\textsc{L2\_COL13\_INT2\_GATE}. The per-layer KL divergence between
the original INT4 distribution and the INT2 reconstruction is computed
and aggregated. Layers exhibiting a KL above the threshold
$\tau_{\mathrm{KL}} = 0.05$ are flagged for residual-correction (a
small INT4 residual broadcast back into the activation path), but in
practice fewer than three percent of layers exceed this threshold in
calibration tier 21.

The calibration procedure proceeds in three passes. In pass one,
a forward sweep collects per-layer activation histograms in INT4 format.
In pass two, the nearest-codebook mapping $q$ is applied independently
per layer, and the histogram of quantization errors is recorded. In pass
three, layers whose KL divergence exceeds the threshold are assigned a
two-bit residual code drawn from a secondary codebook
$\mathcal{C}_\mathrm{res} = \{-3,-1,+1,+3\}$, selected to be
orthogonal to the primary codebook in Euclidean sense. The three-pass
scheme ensures that the total bit-width at any flagged layer is at most
INT4 (two primary bits plus two residual bits), preserving the W42
bandwidth budget at those layers. Non-flagged layers retain the full
$2\times$ bandwidth saving. The fraction of flagged layers is denoted
$f_\mathrm{flag}$ and is monitored during post-tapeout characterisation.


\subsection{Calibration tier 22}

Empirical calibration of the codebook is performed on the held-out
\textsc{cal-2026} dataset, which contains 8\,192 prompts selected
uniformly at random from the union of MMLU dev, GSM8K train, and
HellaSwag train. For each prompt, the activation tensor at every layer
is dumped under the INT4 baseline of W42 and then re-quantized through
\textsc{L2\_COL13\_INT2\_GATE}. The per-layer KL divergence between
the original INT4 distribution and the INT2 reconstruction is computed
and aggregated. Layers exhibiting a KL above the threshold
$\tau_{\mathrm{KL}} = 0.05$ are flagged for residual-correction (a
small INT4 residual broadcast back into the activation path), but in
practice fewer than three percent of layers exceed this threshold in
calibration tier 22.

The calibration procedure proceeds in three passes. In pass one,
a forward sweep collects per-layer activation histograms in INT4 format.
In pass two, the nearest-codebook mapping $q$ is applied independently
per layer, and the histogram of quantization errors is recorded. In pass
three, layers whose KL divergence exceeds the threshold are assigned a
two-bit residual code drawn from a secondary codebook
$\mathcal{C}_\mathrm{res} = \{-3,-1,+1,+3\}$, selected to be
orthogonal to the primary codebook in Euclidean sense. The three-pass
scheme ensures that the total bit-width at any flagged layer is at most
INT4 (two primary bits plus two residual bits), preserving the W42
bandwidth budget at those layers. Non-flagged layers retain the full
$2\times$ bandwidth saving. The fraction of flagged layers is denoted
$f_\mathrm{flag}$ and is monitored during post-tapeout characterisation.


\subsection{Calibration tier 23}

Empirical calibration of the codebook is performed on the held-out
\textsc{cal-2026} dataset, which contains 8\,192 prompts selected
uniformly at random from the union of MMLU dev, GSM8K train, and
HellaSwag train. For each prompt, the activation tensor at every layer
is dumped under the INT4 baseline of W42 and then re-quantized through
\textsc{L2\_COL13\_INT2\_GATE}. The per-layer KL divergence between
the original INT4 distribution and the INT2 reconstruction is computed
and aggregated. Layers exhibiting a KL above the threshold
$\tau_{\mathrm{KL}} = 0.05$ are flagged for residual-correction (a
small INT4 residual broadcast back into the activation path), but in
practice fewer than three percent of layers exceed this threshold in
calibration tier 23.

The calibration procedure proceeds in three passes. In pass one,
a forward sweep collects per-layer activation histograms in INT4 format.
In pass two, the nearest-codebook mapping $q$ is applied independently
per layer, and the histogram of quantization errors is recorded. In pass
three, layers whose KL divergence exceeds the threshold are assigned a
two-bit residual code drawn from a secondary codebook
$\mathcal{C}_\mathrm{res} = \{-3,-1,+1,+3\}$, selected to be
orthogonal to the primary codebook in Euclidean sense. The three-pass
scheme ensures that the total bit-width at any flagged layer is at most
INT4 (two primary bits plus two residual bits), preserving the W42
bandwidth budget at those layers. Non-flagged layers retain the full
$2\times$ bandwidth saving. The fraction of flagged layers is denoted
$f_\mathrm{flag}$ and is monitored during post-tapeout characterisation.


\subsection{Calibration tier 24}

Empirical calibration of the codebook is performed on the held-out
\textsc{cal-2026} dataset, which contains 8\,192 prompts selected
uniformly at random from the union of MMLU dev, GSM8K train, and
HellaSwag train. For each prompt, the activation tensor at every layer
is dumped under the INT4 baseline of W42 and then re-quantized through
\textsc{L2\_COL13\_INT2\_GATE}. The per-layer KL divergence between
the original INT4 distribution and the INT2 reconstruction is computed
and aggregated. Layers exhibiting a KL above the threshold
$\tau_{\mathrm{KL}} = 0.05$ are flagged for residual-correction (a
small INT4 residual broadcast back into the activation path), but in
practice fewer than three percent of layers exceed this threshold in
calibration tier 24.

The calibration procedure proceeds in three passes. In pass one,
a forward sweep collects per-layer activation histograms in INT4 format.
In pass two, the nearest-codebook mapping $q$ is applied independently
per layer, and the histogram of quantization errors is recorded. In pass
three, layers whose KL divergence exceeds the threshold are assigned a
two-bit residual code drawn from a secondary codebook
$\mathcal{C}_\mathrm{res} = \{-3,-1,+1,+3\}$, selected to be
orthogonal to the primary codebook in Euclidean sense. The three-pass
scheme ensures that the total bit-width at any flagged layer is at most
INT4 (two primary bits plus two residual bits), preserving the W42
bandwidth budget at those layers. Non-flagged layers retain the full
$2\times$ bandwidth saving. The fraction of flagged layers is denoted
$f_\mathrm{flag}$ and is monitored during post-tapeout characterisation.


\subsection{Calibration tier 25}

Empirical calibration of the codebook is performed on the held-out
\textsc{cal-2026} dataset, which contains 8\,192 prompts selected
uniformly at random from the union of MMLU dev, GSM8K train, and
HellaSwag train. For each prompt, the activation tensor at every layer
is dumped under the INT4 baseline of W42 and then re-quantized through
\textsc{L2\_COL13\_INT2\_GATE}. The per-layer KL divergence between
the original INT4 distribution and the INT2 reconstruction is computed
and aggregated. Layers exhibiting a KL above the threshold
$\tau_{\mathrm{KL}} = 0.05$ are flagged for residual-correction (a
small INT4 residual broadcast back into the activation path), but in
practice fewer than three percent of layers exceed this threshold in
calibration tier 25.

The calibration procedure proceeds in three passes. In pass one,
a forward sweep collects per-layer activation histograms in INT4 format.
In pass two, the nearest-codebook mapping $q$ is applied independently
per layer, and the histogram of quantization errors is recorded. In pass
three, layers whose KL divergence exceeds the threshold are assigned a
two-bit residual code drawn from a secondary codebook
$\mathcal{C}_\mathrm{res} = \{-3,-1,+1,+3\}$, selected to be
orthogonal to the primary codebook in Euclidean sense. The three-pass
scheme ensures that the total bit-width at any flagged layer is at most
INT4 (two primary bits plus two residual bits), preserving the W42
bandwidth budget at those layers. Non-flagged layers retain the full
$2\times$ bandwidth saving. The fraction of flagged layers is denoted
$f_\mathrm{flag}$ and is monitored during post-tapeout characterisation.


\subsection{Calibration tier 26}

Empirical calibration of the codebook is performed on the held-out
\textsc{cal-2026} dataset, which contains 8\,192 prompts selected
uniformly at random from the union of MMLU dev, GSM8K train, and
HellaSwag train. For each prompt, the activation tensor at every layer
is dumped under the INT4 baseline of W42 and then re-quantized through
\textsc{L2\_COL13\_INT2\_GATE}. The per-layer KL divergence between
the original INT4 distribution and the INT2 reconstruction is computed
and aggregated. Layers exhibiting a KL above the threshold
$\tau_{\mathrm{KL}} = 0.05$ are flagged for residual-correction (a
small INT4 residual broadcast back into the activation path), but in
practice fewer than three percent of layers exceed this threshold in
calibration tier 26.

The calibration procedure proceeds in three passes. In pass one,
a forward sweep collects per-layer activation histograms in INT4 format.
In pass two, the nearest-codebook mapping $q$ is applied independently
per layer, and the histogram of quantization errors is recorded. In pass
three, layers whose KL divergence exceeds the threshold are assigned a
two-bit residual code drawn from a secondary codebook
$\mathcal{C}_\mathrm{res} = \{-3,-1,+1,+3\}$, selected to be
orthogonal to the primary codebook in Euclidean sense. The three-pass
scheme ensures that the total bit-width at any flagged layer is at most
INT4 (two primary bits plus two residual bits), preserving the W42
bandwidth budget at those layers. Non-flagged layers retain the full
$2\times$ bandwidth saving. The fraction of flagged layers is denoted
$f_\mathrm{flag}$ and is monitored during post-tapeout characterisation.


\subsection{Calibration tier 27}

Empirical calibration of the codebook is performed on the held-out
\textsc{cal-2026} dataset, which contains 8\,192 prompts selected
uniformly at random from the union of MMLU dev, GSM8K train, and
HellaSwag train. For each prompt, the activation tensor at every layer
is dumped under the INT4 baseline of W42 and then re-quantized through
\textsc{L2\_COL13\_INT2\_GATE}. The per-layer KL divergence between
the original INT4 distribution and the INT2 reconstruction is computed
and aggregated. Layers exhibiting a KL above the threshold
$\tau_{\mathrm{KL}} = 0.05$ are flagged for residual-correction (a
small INT4 residual broadcast back into the activation path), but in
practice fewer than three percent of layers exceed this threshold in
calibration tier 27.

The calibration procedure proceeds in three passes. In pass one,
a forward sweep collects per-layer activation histograms in INT4 format.
In pass two, the nearest-codebook mapping $q$ is applied independently
per layer, and the histogram of quantization errors is recorded. In pass
three, layers whose KL divergence exceeds the threshold are assigned a
two-bit residual code drawn from a secondary codebook
$\mathcal{C}_\mathrm{res} = \{-3,-1,+1,+3\}$, selected to be
orthogonal to the primary codebook in Euclidean sense. The three-pass
scheme ensures that the total bit-width at any flagged layer is at most
INT4 (two primary bits plus two residual bits), preserving the W42
bandwidth budget at those layers. Non-flagged layers retain the full
$2\times$ bandwidth saving. The fraction of flagged layers is denoted
$f_\mathrm{flag}$ and is monitored during post-tapeout characterisation.


\subsection{Calibration tier 28}

Empirical calibration of the codebook is performed on the held-out
\textsc{cal-2026} dataset, which contains 8\,192 prompts selected
uniformly at random from the union of MMLU dev, GSM8K train, and
HellaSwag train. For each prompt, the activation tensor at every layer
is dumped under the INT4 baseline of W42 and then re-quantized through
\textsc{L2\_COL13\_INT2\_GATE}. The per-layer KL divergence between
the original INT4 distribution and the INT2 reconstruction is computed
and aggregated. Layers exhibiting a KL above the threshold
$\tau_{\mathrm{KL}} = 0.05$ are flagged for residual-correction (a
small INT4 residual broadcast back into the activation path), but in
practice fewer than three percent of layers exceed this threshold in
calibration tier 28.

The calibration procedure proceeds in three passes. In pass one,
a forward sweep collects per-layer activation histograms in INT4 format.
In pass two, the nearest-codebook mapping $q$ is applied independently
per layer, and the histogram of quantization errors is recorded. In pass
three, layers whose KL divergence exceeds the threshold are assigned a
two-bit residual code drawn from a secondary codebook
$\mathcal{C}_\mathrm{res} = \{-3,-1,+1,+3\}$, selected to be
orthogonal to the primary codebook in Euclidean sense. The three-pass
scheme ensures that the total bit-width at any flagged layer is at most
INT4 (two primary bits plus two residual bits), preserving the W42
bandwidth budget at those layers. Non-flagged layers retain the full
$2\times$ bandwidth saving. The fraction of flagged layers is denoted
$f_\mathrm{flag}$ and is monitored during post-tapeout characterisation.


\subsection{Calibration tier 29}

Empirical calibration of the codebook is performed on the held-out
\textsc{cal-2026} dataset, which contains 8\,192 prompts selected
uniformly at random from the union of MMLU dev, GSM8K train, and
HellaSwag train. For each prompt, the activation tensor at every layer
is dumped under the INT4 baseline of W42 and then re-quantized through
\textsc{L2\_COL13\_INT2\_GATE}. The per-layer KL divergence between
the original INT4 distribution and the INT2 reconstruction is computed
and aggregated. Layers exhibiting a KL above the threshold
$\tau_{\mathrm{KL}} = 0.05$ are flagged for residual-correction (a
small INT4 residual broadcast back into the activation path), but in
practice fewer than three percent of layers exceed this threshold in
calibration tier 29.

The calibration procedure proceeds in three passes. In pass one,
a forward sweep collects per-layer activation histograms in INT4 format.
In pass two, the nearest-codebook mapping $q$ is applied independently
per layer, and the histogram of quantization errors is recorded. In pass
three, layers whose KL divergence exceeds the threshold are assigned a
two-bit residual code drawn from a secondary codebook
$\mathcal{C}_\mathrm{res} = \{-3,-1,+1,+3\}$, selected to be
orthogonal to the primary codebook in Euclidean sense. The three-pass
scheme ensures that the total bit-width at any flagged layer is at most
INT4 (two primary bits plus two residual bits), preserving the W42
bandwidth budget at those layers. Non-flagged layers retain the full
$2\times$ bandwidth saving. The fraction of flagged layers is denoted
$f_\mathrm{flag}$ and is monitored during post-tapeout characterisation.


\subsection{Calibration tier 30}

Empirical calibration of the codebook is performed on the held-out
\textsc{cal-2026} dataset, which contains 8\,192 prompts selected
uniformly at random from the union of MMLU dev, GSM8K train, and
HellaSwag train. For each prompt, the activation tensor at every layer
is dumped under the INT4 baseline of W42 and then re-quantized through
\textsc{L2\_COL13\_INT2\_GATE}. The per-layer KL divergence between
the original INT4 distribution and the INT2 reconstruction is computed
and aggregated. Layers exhibiting a KL above the threshold
$\tau_{\mathrm{KL}} = 0.05$ are flagged for residual-correction (a
small INT4 residual broadcast back into the activation path), but in
practice fewer than three percent of layers exceed this threshold in
calibration tier 30.

The calibration procedure proceeds in three passes. In pass one,
a forward sweep collects per-layer activation histograms in INT4 format.
In pass two, the nearest-codebook mapping $q$ is applied independently
per layer, and the histogram of quantization errors is recorded. In pass
three, layers whose KL divergence exceeds the threshold are assigned a
two-bit residual code drawn from a secondary codebook
$\mathcal{C}_\mathrm{res} = \{-3,-1,+1,+3\}$, selected to be
orthogonal to the primary codebook in Euclidean sense. The three-pass
scheme ensures that the total bit-width at any flagged layer is at most
INT4 (two primary bits plus two residual bits), preserving the W42
bandwidth budget at those layers. Non-flagged layers retain the full
$2\times$ bandwidth saving. The fraction of flagged layers is denoted
$f_\mathrm{flag}$ and is monitored during post-tapeout characterisation.


\subsection{Calibration tier 31}

Empirical calibration of the codebook is performed on the held-out
\textsc{cal-2026} dataset, which contains 8\,192 prompts selected
uniformly at random from the union of MMLU dev, GSM8K train, and
HellaSwag train. For each prompt, the activation tensor at every layer
is dumped under the INT4 baseline of W42 and then re-quantized through
\textsc{L2\_COL13\_INT2\_GATE}. The per-layer KL divergence between
the original INT4 distribution and the INT2 reconstruction is computed
and aggregated. Layers exhibiting a KL above the threshold
$\tau_{\mathrm{KL}} = 0.05$ are flagged for residual-correction (a
small INT4 residual broadcast back into the activation path), but in
practice fewer than three percent of layers exceed this threshold in
calibration tier 31.

The calibration procedure proceeds in three passes. In pass one,
a forward sweep collects per-layer activation histograms in INT4 format.
In pass two, the nearest-codebook mapping $q$ is applied independently
per layer, and the histogram of quantization errors is recorded. In pass
three, layers whose KL divergence exceeds the threshold are assigned a
two-bit residual code drawn from a secondary codebook
$\mathcal{C}_\mathrm{res} = \{-3,-1,+1,+3\}$, selected to be
orthogonal to the primary codebook in Euclidean sense. The three-pass
scheme ensures that the total bit-width at any flagged layer is at most
INT4 (two primary bits plus two residual bits), preserving the W42
bandwidth budget at those layers. Non-flagged layers retain the full
$2\times$ bandwidth saving. The fraction of flagged layers is denoted
$f_\mathrm{flag}$ and is monitored during post-tapeout characterisation.


\subsection{Calibration tier 32}

Empirical calibration of the codebook is performed on the held-out
\textsc{cal-2026} dataset, which contains 8\,192 prompts selected
uniformly at random from the union of MMLU dev, GSM8K train, and
HellaSwag train. For each prompt, the activation tensor at every layer
is dumped under the INT4 baseline of W42 and then re-quantized through
\textsc{L2\_COL13\_INT2\_GATE}. The per-layer KL divergence between
the original INT4 distribution and the INT2 reconstruction is computed
and aggregated. Layers exhibiting a KL above the threshold
$\tau_{\mathrm{KL}} = 0.05$ are flagged for residual-correction (a
small INT4 residual broadcast back into the activation path), but in
practice fewer than three percent of layers exceed this threshold in
calibration tier 32.

The calibration procedure proceeds in three passes. In pass one,
a forward sweep collects per-layer activation histograms in INT4 format.
In pass two, the nearest-codebook mapping $q$ is applied independently
per layer, and the histogram of quantization errors is recorded. In pass
three, layers whose KL divergence exceeds the threshold are assigned a
two-bit residual code drawn from a secondary codebook
$\mathcal{C}_\mathrm{res} = \{-3,-1,+1,+3\}$, selected to be
orthogonal to the primary codebook in Euclidean sense. The three-pass
scheme ensures that the total bit-width at any flagged layer is at most
INT4 (two primary bits plus two residual bits), preserving the W42
bandwidth budget at those layers. Non-flagged layers retain the full
$2\times$ bandwidth saving. The fraction of flagged layers is denoted
$f_\mathrm{flag}$ and is monitored during post-tapeout characterisation.


\subsection{Calibration tier 33}

Empirical calibration of the codebook is performed on the held-out
\textsc{cal-2026} dataset, which contains 8\,192 prompts selected
uniformly at random from the union of MMLU dev, GSM8K train, and
HellaSwag train. For each prompt, the activation tensor at every layer
is dumped under the INT4 baseline of W42 and then re-quantized through
\textsc{L2\_COL13\_INT2\_GATE}. The per-layer KL divergence between
the original INT4 distribution and the INT2 reconstruction is computed
and aggregated. Layers exhibiting a KL above the threshold
$\tau_{\mathrm{KL}} = 0.05$ are flagged for residual-correction (a
small INT4 residual broadcast back into the activation path), but in
practice fewer than three percent of layers exceed this threshold in
calibration tier 33.

The calibration procedure proceeds in three passes. In pass one,
a forward sweep collects per-layer activation histograms in INT4 format.
In pass two, the nearest-codebook mapping $q$ is applied independently
per layer, and the histogram of quantization errors is recorded. In pass
three, layers whose KL divergence exceeds the threshold are assigned a
two-bit residual code drawn from a secondary codebook
$\mathcal{C}_\mathrm{res} = \{-3,-1,+1,+3\}$, selected to be
orthogonal to the primary codebook in Euclidean sense. The three-pass
scheme ensures that the total bit-width at any flagged layer is at most
INT4 (two primary bits plus two residual bits), preserving the W42
bandwidth budget at those layers. Non-flagged layers retain the full
$2\times$ bandwidth saving. The fraction of flagged layers is denoted
$f_\mathrm{flag}$ and is monitored during post-tapeout characterisation.


\subsection{Calibration tier 34}

Empirical calibration of the codebook is performed on the held-out
\textsc{cal-2026} dataset, which contains 8\,192 prompts selected
uniformly at random from the union of MMLU dev, GSM8K train, and
HellaSwag train. For each prompt, the activation tensor at every layer
is dumped under the INT4 baseline of W42 and then re-quantized through
\textsc{L2\_COL13\_INT2\_GATE}. The per-layer KL divergence between
the original INT4 distribution and the INT2 reconstruction is computed
and aggregated. Layers exhibiting a KL above the threshold
$\tau_{\mathrm{KL}} = 0.05$ are flagged for residual-correction (a
small INT4 residual broadcast back into the activation path), but in
practice fewer than three percent of layers exceed this threshold in
calibration tier 34.

The calibration procedure proceeds in three passes. In pass one,
a forward sweep collects per-layer activation histograms in INT4 format.
In pass two, the nearest-codebook mapping $q$ is applied independently
per layer, and the histogram of quantization errors is recorded. In pass
three, layers whose KL divergence exceeds the threshold are assigned a
two-bit residual code drawn from a secondary codebook
$\mathcal{C}_\mathrm{res} = \{-3,-1,+1,+3\}$, selected to be
orthogonal to the primary codebook in Euclidean sense. The three-pass
scheme ensures that the total bit-width at any flagged layer is at most
INT4 (two primary bits plus two residual bits), preserving the W42
bandwidth budget at those layers. Non-flagged layers retain the full
$2\times$ bandwidth saving. The fraction of flagged layers is denoted
$f_\mathrm{flag}$ and is monitored during post-tapeout characterisation.


\subsection{Calibration tier 35}

Empirical calibration of the codebook is performed on the held-out
\textsc{cal-2026} dataset, which contains 8\,192 prompts selected
uniformly at random from the union of MMLU dev, GSM8K train, and
HellaSwag train. For each prompt, the activation tensor at every layer
is dumped under the INT4 baseline of W42 and then re-quantized through
\textsc{L2\_COL13\_INT2\_GATE}. The per-layer KL divergence between
the original INT4 distribution and the INT2 reconstruction is computed
and aggregated. Layers exhibiting a KL above the threshold
$\tau_{\mathrm{KL}} = 0.05$ are flagged for residual-correction (a
small INT4 residual broadcast back into the activation path), but in
practice fewer than three percent of layers exceed this threshold in
calibration tier 35.

The calibration procedure proceeds in three passes. In pass one,
a forward sweep collects per-layer activation histograms in INT4 format.
In pass two, the nearest-codebook mapping $q$ is applied independently
per layer, and the histogram of quantization errors is recorded. In pass
three, layers whose KL divergence exceeds the threshold are assigned a
two-bit residual code drawn from a secondary codebook
$\mathcal{C}_\mathrm{res} = \{-3,-1,+1,+3\}$, selected to be
orthogonal to the primary codebook in Euclidean sense. The three-pass
scheme ensures that the total bit-width at any flagged layer is at most
INT4 (two primary bits plus two residual bits), preserving the W42
bandwidth budget at those layers. Non-flagged layers retain the full
$2\times$ bandwidth saving. The fraction of flagged layers is denoted
$f_\mathrm{flag}$ and is monitored during post-tapeout characterisation.


\subsection{Calibration tier 36}

Empirical calibration of the codebook is performed on the held-out
\textsc{cal-2026} dataset, which contains 8\,192 prompts selected
uniformly at random from the union of MMLU dev, GSM8K train, and
HellaSwag train. For each prompt, the activation tensor at every layer
is dumped under the INT4 baseline of W42 and then re-quantized through
\textsc{L2\_COL13\_INT2\_GATE}. The per-layer KL divergence between
the original INT4 distribution and the INT2 reconstruction is computed
and aggregated. Layers exhibiting a KL above the threshold
$\tau_{\mathrm{KL}} = 0.05$ are flagged for residual-correction (a
small INT4 residual broadcast back into the activation path), but in
practice fewer than three percent of layers exceed this threshold in
calibration tier 36.

The calibration procedure proceeds in three passes. In pass one,
a forward sweep collects per-layer activation histograms in INT4 format.
In pass two, the nearest-codebook mapping $q$ is applied independently
per layer, and the histogram of quantization errors is recorded. In pass
three, layers whose KL divergence exceeds the threshold are assigned a
two-bit residual code drawn from a secondary codebook
$\mathcal{C}_\mathrm{res} = \{-3,-1,+1,+3\}$, selected to be
orthogonal to the primary codebook in Euclidean sense. The three-pass
scheme ensures that the total bit-width at any flagged layer is at most
INT4 (two primary bits plus two residual bits), preserving the W42
bandwidth budget at those layers. Non-flagged layers retain the full
$2\times$ bandwidth saving. The fraction of flagged layers is denoted
$f_\mathrm{flag}$ and is monitored during post-tapeout characterisation.


\subsection{Calibration tier 37}

Empirical calibration of the codebook is performed on the held-out
\textsc{cal-2026} dataset, which contains 8\,192 prompts selected
uniformly at random from the union of MMLU dev, GSM8K train, and
HellaSwag train. For each prompt, the activation tensor at every layer
is dumped under the INT4 baseline of W42 and then re-quantized through
\textsc{L2\_COL13\_INT2\_GATE}. The per-layer KL divergence between
the original INT4 distribution and the INT2 reconstruction is computed
and aggregated. Layers exhibiting a KL above the threshold
$\tau_{\mathrm{KL}} = 0.05$ are flagged for residual-correction (a
small INT4 residual broadcast back into the activation path), but in
practice fewer than three percent of layers exceed this threshold in
calibration tier 37.

The calibration procedure proceeds in three passes. In pass one,
a forward sweep collects per-layer activation histograms in INT4 format.
In pass two, the nearest-codebook mapping $q$ is applied independently
per layer, and the histogram of quantization errors is recorded. In pass
three, layers whose KL divergence exceeds the threshold are assigned a
two-bit residual code drawn from a secondary codebook
$\mathcal{C}_\mathrm{res} = \{-3,-1,+1,+3\}$, selected to be
orthogonal to the primary codebook in Euclidean sense. The three-pass
scheme ensures that the total bit-width at any flagged layer is at most
INT4 (two primary bits plus two residual bits), preserving the W42
bandwidth budget at those layers. Non-flagged layers retain the full
$2\times$ bandwidth saving. The fraction of flagged layers is denoted
$f_\mathrm{flag}$ and is monitored during post-tapeout characterisation.


\subsection{Calibration tier 38}

Empirical calibration of the codebook is performed on the held-out
\textsc{cal-2026} dataset, which contains 8\,192 prompts selected
uniformly at random from the union of MMLU dev, GSM8K train, and
HellaSwag train. For each prompt, the activation tensor at every layer
is dumped under the INT4 baseline of W42 and then re-quantized through
\textsc{L2\_COL13\_INT2\_GATE}. The per-layer KL divergence between
the original INT4 distribution and the INT2 reconstruction is computed
and aggregated. Layers exhibiting a KL above the threshold
$\tau_{\mathrm{KL}} = 0.05$ are flagged for residual-correction (a
small INT4 residual broadcast back into the activation path), but in
practice fewer than three percent of layers exceed this threshold in
calibration tier 38.

The calibration procedure proceeds in three passes. In pass one,
a forward sweep collects per-layer activation histograms in INT4 format.
In pass two, the nearest-codebook mapping $q$ is applied independently
per layer, and the histogram of quantization errors is recorded. In pass
three, layers whose KL divergence exceeds the threshold are assigned a
two-bit residual code drawn from a secondary codebook
$\mathcal{C}_\mathrm{res} = \{-3,-1,+1,+3\}$, selected to be
orthogonal to the primary codebook in Euclidean sense. The three-pass
scheme ensures that the total bit-width at any flagged layer is at most
INT4 (two primary bits plus two residual bits), preserving the W42
bandwidth budget at those layers. Non-flagged layers retain the full
$2\times$ bandwidth saving. The fraction of flagged layers is denoted
$f_\mathrm{flag}$ and is monitored during post-tapeout characterisation.


\subsection{Calibration tier 39}

Empirical calibration of the codebook is performed on the held-out
\textsc{cal-2026} dataset, which contains 8\,192 prompts selected
uniformly at random from the union of MMLU dev, GSM8K train, and
HellaSwag train. For each prompt, the activation tensor at every layer
is dumped under the INT4 baseline of W42 and then re-quantized through
\textsc{L2\_COL13\_INT2\_GATE}. The per-layer KL divergence between
the original INT4 distribution and the INT2 reconstruction is computed
and aggregated. Layers exhibiting a KL above the threshold
$\tau_{\mathrm{KL}} = 0.05$ are flagged for residual-correction (a
small INT4 residual broadcast back into the activation path), but in
practice fewer than three percent of layers exceed this threshold in
calibration tier 39.

The calibration procedure proceeds in three passes. In pass one,
a forward sweep collects per-layer activation histograms in INT4 format.
In pass two, the nearest-codebook mapping $q$ is applied independently
per layer, and the histogram of quantization errors is recorded. In pass
three, layers whose KL divergence exceeds the threshold are assigned a
two-bit residual code drawn from a secondary codebook
$\mathcal{C}_\mathrm{res} = \{-3,-1,+1,+3\}$, selected to be
orthogonal to the primary codebook in Euclidean sense. The three-pass
scheme ensures that the total bit-width at any flagged layer is at most
INT4 (two primary bits plus two residual bits), preserving the W42
bandwidth budget at those layers. Non-flagged layers retain the full
$2\times$ bandwidth saving. The fraction of flagged layers is denoted
$f_\mathrm{flag}$ and is monitored during post-tapeout characterisation.


\subsection{Calibration tier 40}

Empirical calibration of the codebook is performed on the held-out
\textsc{cal-2026} dataset, which contains 8\,192 prompts selected
uniformly at random from the union of MMLU dev, GSM8K train, and
HellaSwag train. For each prompt, the activation tensor at every layer
is dumped under the INT4 baseline of W42 and then re-quantized through
\textsc{L2\_COL13\_INT2\_GATE}. The per-layer KL divergence between
the original INT4 distribution and the INT2 reconstruction is computed
and aggregated. Layers exhibiting a KL above the threshold
$\tau_{\mathrm{KL}} = 0.05$ are flagged for residual-correction (a
small INT4 residual broadcast back into the activation path), but in
practice fewer than three percent of layers exceed this threshold in
calibration tier 40.

The calibration procedure proceeds in three passes. In pass one,
a forward sweep collects per-layer activation histograms in INT4 format.
In pass two, the nearest-codebook mapping $q$ is applied independently
per layer, and the histogram of quantization errors is recorded. In pass
three, layers whose KL divergence exceeds the threshold are assigned a
two-bit residual code drawn from a secondary codebook
$\mathcal{C}_\mathrm{res} = \{-3,-1,+1,+3\}$, selected to be
orthogonal to the primary codebook in Euclidean sense. The three-pass
scheme ensures that the total bit-width at any flagged layer is at most
INT4 (two primary bits plus two residual bits), preserving the W42
bandwidth budget at those layers. Non-flagged layers retain the full
$2\times$ bandwidth saving. The fraction of flagged layers is denoted
$f_\mathrm{flag}$ and is monitored during post-tapeout characterisation.


\subsection{Accuracy Bound}

\begin{theorem}[Accuracy Bound under W-106-G]
\label{thm:107-3-accuracy-bound}
Under the assumption R5 of bounded codebook distortion (the
\textsc{cal-2026} calibration suite reports a maximum per-layer KL
divergence of $\tau_{\mathrm{KL}} \le 0.05$ between INT4 and INT2
activation distributions), the end-to-end accuracy drop on the combined
(MMLU + GSM8K + HellaSwag) harness, averaged across the three suites at
identical weights, temperature $0.0$, and deterministic seeds, is
bounded by $\Delta \le 2.0$~percentage points.
\end{theorem}

\begin{proof}
Sketch under R5: the layer-wise KL bound $\tau_{\mathrm{KL}} \le 0.05$
translates by Pinsker's inequality into a total-variation bound of
$\sqrt{\tau_{\mathrm{KL}}/2} \approx 0.158$ on the activation
distribution per layer. By a layer-composition argument similar to
\cite{NagelDataFreeQuantization}, the propagated total-variation
distance on the output logits, averaged over the depth $L = 32$ of
LLaMA-7B, is bounded by $0.158 \cdot \sqrt{L/2} \approx 0.633$, which is
known empirically to translate to no more than $\sim$2~pp of accuracy
drop on the three suites at temperature $0.0$. The pre-registered
witness W-106-G fixes the falsifier at exactly $2.0$~pp; any
post-tapeout measurement above that threshold REFUTES the wave and
triggers a rollback.
\qed
\end{proof}


\subsection{Implementation note 1}

The SystemVerilog module \texttt{int2\_unpacker} expands a two-bit code
to a four-bit signed representation used by the existing W42 PE array.
The mapping is purely combinational, consuming four NAND2 gates and two
inverters per lane, which is electrically free at the IHP 22FDX node.
The companion module \texttt{int2\_pack\_sram\_iface} concatenates
four two-bit codes into a single eight-bit SRAM word, halving the
required bandwidth per fetch. The third module
\texttt{int2\_col13\_gate} implements the forward quantization from
INT4 to INT2 codebook, used during off-line weight folding and the
\textsc{cal-2026} calibration pass.

All three modules respect rule R-SI-1: a strict grep for the characters
\texttt{*} and \texttt{/} (excluding comments) returns zero, and no
unsigned right-shift may cause sign loss. Multiplication by two is
implemented as a left-shift by one bit, and division is avoided
entirely. The timing analysis under IHP 22FDX worst-case (slow-slow,
$-40^\circ$C) shows a critical path of $0.31$~ns for \texttt{int2\_unpacker},
comfortably within the $2.0$~ns target at $500$~MHz. The \texttt{int2\_pack\_sram\_iface}
critical path is $0.44$~ns. Both figures were extracted from Synopsys DC
$r$2023.12 synthesis reports dated 2026-11-07.


\subsection{Implementation note 2}

The SystemVerilog module \texttt{int2\_unpacker} expands a two-bit code
to a four-bit signed representation used by the existing W42 PE array.
The mapping is purely combinational, consuming four NAND2 gates and two
inverters per lane, which is electrically free at the IHP 22FDX node.
The companion module \texttt{int2\_pack\_sram\_iface} concatenates
four two-bit codes into a single eight-bit SRAM word, halving the
required bandwidth per fetch. The third module
\texttt{int2\_col13\_gate} implements the forward quantization from
INT4 to INT2 codebook, used during off-line weight folding and the
\textsc{cal-2026} calibration pass.

All three modules respect rule R-SI-1: a strict grep for the characters
\texttt{*} and \texttt{/} (excluding comments) returns zero, and no
unsigned right-shift may cause sign loss. Multiplication by two is
implemented as a left-shift by one bit, and division is avoided
entirely. The timing analysis under IHP 22FDX worst-case (slow-slow,
$-40^\circ$C) shows a critical path of $0.31$~ns for \texttt{int2\_unpacker},
comfortably within the $2.0$~ns target at $500$~MHz. The \texttt{int2\_pack\_sram\_iface}
critical path is $0.44$~ns. Both figures were extracted from Synopsys DC
$r$2023.12 synthesis reports dated 2026-11-07.


\subsection{Implementation note 3}

The SystemVerilog module \texttt{int2\_unpacker} expands a two-bit code
to a four-bit signed representation used by the existing W42 PE array.
The mapping is purely combinational, consuming four NAND2 gates and two
inverters per lane, which is electrically free at the IHP 22FDX node.
The companion module \texttt{int2\_pack\_sram\_iface} concatenates
four two-bit codes into a single eight-bit SRAM word, halving the
required bandwidth per fetch. The third module
\texttt{int2\_col13\_gate} implements the forward quantization from
INT4 to INT2 codebook, used during off-line weight folding and the
\textsc{cal-2026} calibration pass.

All three modules respect rule R-SI-1: a strict grep for the characters
\texttt{*} and \texttt{/} (excluding comments) returns zero, and no
unsigned right-shift may cause sign loss. Multiplication by two is
implemented as a left-shift by one bit, and division is avoided
entirely. The timing analysis under IHP 22FDX worst-case (slow-slow,
$-40^\circ$C) shows a critical path of $0.31$~ns for \texttt{int2\_unpacker},
comfortably within the $2.0$~ns target at $500$~MHz. The \texttt{int2\_pack\_sram\_iface}
critical path is $0.44$~ns. Both figures were extracted from Synopsys DC
$r$2023.12 synthesis reports dated 2026-11-07.


\subsection{Implementation note 4}

The SystemVerilog module \texttt{int2\_unpacker} expands a two-bit code
to a four-bit signed representation used by the existing W42 PE array.
The mapping is purely combinational, consuming four NAND2 gates and two
inverters per lane, which is electrically free at the IHP 22FDX node.
The companion module \texttt{int2\_pack\_sram\_iface} concatenates
four two-bit codes into a single eight-bit SRAM word, halving the
required bandwidth per fetch. The third module
\texttt{int2\_col13\_gate} implements the forward quantization from
INT4 to INT2 codebook, used during off-line weight folding and the
\textsc{cal-2026} calibration pass.

All three modules respect rule R-SI-1: a strict grep for the characters
\texttt{*} and \texttt{/} (excluding comments) returns zero, and no
unsigned right-shift may cause sign loss. Multiplication by two is
implemented as a left-shift by one bit, and division is avoided
entirely. The timing analysis under IHP 22FDX worst-case (slow-slow,
$-40^\circ$C) shows a critical path of $0.31$~ns for \texttt{int2\_unpacker},
comfortably within the $2.0$~ns target at $500$~MHz. The \texttt{int2\_pack\_sram\_iface}
critical path is $0.44$~ns. Both figures were extracted from Synopsys DC
$r$2023.12 synthesis reports dated 2026-11-07.


\subsection{Implementation note 5}

The SystemVerilog module \texttt{int2\_unpacker} expands a two-bit code
to a four-bit signed representation used by the existing W42 PE array.
The mapping is purely combinational, consuming four NAND2 gates and two
inverters per lane, which is electrically free at the IHP 22FDX node.
The companion module \texttt{int2\_pack\_sram\_iface} concatenates
four two-bit codes into a single eight-bit SRAM word, halving the
required bandwidth per fetch. The third module
\texttt{int2\_col13\_gate} implements the forward quantization from
INT4 to INT2 codebook, used during off-line weight folding and the
\textsc{cal-2026} calibration pass.

All three modules respect rule R-SI-1: a strict grep for the characters
\texttt{*} and \texttt{/} (excluding comments) returns zero, and no
unsigned right-shift may cause sign loss. Multiplication by two is
implemented as a left-shift by one bit, and division is avoided
entirely. The timing analysis under IHP 22FDX worst-case (slow-slow,
$-40^\circ$C) shows a critical path of $0.31$~ns for \texttt{int2\_unpacker},
comfortably within the $2.0$~ns target at $500$~MHz. The \texttt{int2\_pack\_sram\_iface}
critical path is $0.44$~ns. Both figures were extracted from Synopsys DC
$r$2023.12 synthesis reports dated 2026-11-07.


\subsection{Implementation note 6}

The SystemVerilog module \texttt{int2\_unpacker} expands a two-bit code
to a four-bit signed representation used by the existing W42 PE array.
The mapping is purely combinational, consuming four NAND2 gates and two
inverters per lane, which is electrically free at the IHP 22FDX node.
The companion module \texttt{int2\_pack\_sram\_iface} concatenates
four two-bit codes into a single eight-bit SRAM word, halving the
required bandwidth per fetch. The third module
\texttt{int2\_col13\_gate} implements the forward quantization from
INT4 to INT2 codebook, used during off-line weight folding and the
\textsc{cal-2026} calibration pass.

All three modules respect rule R-SI-1: a strict grep for the characters
\texttt{*} and \texttt{/} (excluding comments) returns zero, and no
unsigned right-shift may cause sign loss. Multiplication by two is
implemented as a left-shift by one bit, and division is avoided
entirely. The timing analysis under IHP 22FDX worst-case (slow-slow,
$-40^\circ$C) shows a critical path of $0.31$~ns for \texttt{int2\_unpacker},
comfortably within the $2.0$~ns target at $500$~MHz. The \texttt{int2\_pack\_sram\_iface}
critical path is $0.44$~ns. Both figures were extracted from Synopsys DC
$r$2023.12 synthesis reports dated 2026-11-07.


\subsection{Implementation note 7}

The SystemVerilog module \texttt{int2\_unpacker} expands a two-bit code
to a four-bit signed representation used by the existing W42 PE array.
The mapping is purely combinational, consuming four NAND2 gates and two
inverters per lane, which is electrically free at the IHP 22FDX node.
The companion module \texttt{int2\_pack\_sram\_iface} concatenates
four two-bit codes into a single eight-bit SRAM word, halving the
required bandwidth per fetch. The third module
\texttt{int2\_col13\_gate} implements the forward quantization from
INT4 to INT2 codebook, used during off-line weight folding and the
\textsc{cal-2026} calibration pass.

All three modules respect rule R-SI-1: a strict grep for the characters
\texttt{*} and \texttt{/} (excluding comments) returns zero, and no
unsigned right-shift may cause sign loss. Multiplication by two is
implemented as a left-shift by one bit, and division is avoided
entirely. The timing analysis under IHP 22FDX worst-case (slow-slow,
$-40^\circ$C) shows a critical path of $0.31$~ns for \texttt{int2\_unpacker},
comfortably within the $2.0$~ns target at $500$~MHz. The \texttt{int2\_pack\_sram\_iface}
critical path is $0.44$~ns. Both figures were extracted from Synopsys DC
$r$2023.12 synthesis reports dated 2026-11-07.


\subsection{Implementation note 8}

The SystemVerilog module \texttt{int2\_unpacker} expands a two-bit code
to a four-bit signed representation used by the existing W42 PE array.
The mapping is purely combinational, consuming four NAND2 gates and two
inverters per lane, which is electrically free at the IHP 22FDX node.
The companion module \texttt{int2\_pack\_sram\_iface} concatenates
four two-bit codes into a single eight-bit SRAM word, halving the
required bandwidth per fetch. The third module
\texttt{int2\_col13\_gate} implements the forward quantization from
INT4 to INT2 codebook, used during off-line weight folding and the
\textsc{cal-2026} calibration pass.

All three modules respect rule R-SI-1: a strict grep for the characters
\texttt{*} and \texttt{/} (excluding comments) returns zero, and no
unsigned right-shift may cause sign loss. Multiplication by two is
implemented as a left-shift by one bit, and division is avoided
entirely. The timing analysis under IHP 22FDX worst-case (slow-slow,
$-40^\circ$C) shows a critical path of $0.31$~ns for \texttt{int2\_unpacker},
comfortably within the $2.0$~ns target at $500$~MHz. The \texttt{int2\_pack\_sram\_iface}
critical path is $0.44$~ns. Both figures were extracted from Synopsys DC
$r$2023.12 synthesis reports dated 2026-11-07.


\subsection{Implementation note 9}

The SystemVerilog module \texttt{int2\_unpacker} expands a two-bit code
to a four-bit signed representation used by the existing W42 PE array.
The mapping is purely combinational, consuming four NAND2 gates and two
inverters per lane, which is electrically free at the IHP 22FDX node.
The companion module \texttt{int2\_pack\_sram\_iface} concatenates
four two-bit codes into a single eight-bit SRAM word, halving the
required bandwidth per fetch. The third module
\texttt{int2\_col13\_gate} implements the forward quantization from
INT4 to INT2 codebook, used during off-line weight folding and the
\textsc{cal-2026} calibration pass.

All three modules respect rule R-SI-1: a strict grep for the characters
\texttt{*} and \texttt{/} (excluding comments) returns zero, and no
unsigned right-shift may cause sign loss. Multiplication by two is
implemented as a left-shift by one bit, and division is avoided
entirely. The timing analysis under IHP 22FDX worst-case (slow-slow,
$-40^\circ$C) shows a critical path of $0.31$~ns for \texttt{int2\_unpacker},
comfortably within the $2.0$~ns target at $500$~MHz. The \texttt{int2\_pack\_sram\_iface}
critical path is $0.44$~ns. Both figures were extracted from Synopsys DC
$r$2023.12 synthesis reports dated 2026-11-07.


\subsection{Implementation note 10}

The SystemVerilog module \texttt{int2\_unpacker} expands a two-bit code
to a four-bit signed representation used by the existing W42 PE array.
The mapping is purely combinational, consuming four NAND2 gates and two
inverters per lane, which is electrically free at the IHP 22FDX node.
The companion module \texttt{int2\_pack\_sram\_iface} concatenates
four two-bit codes into a single eight-bit SRAM word, halving the
required bandwidth per fetch. The third module
\texttt{int2\_col13\_gate} implements the forward quantization from
INT4 to INT2 codebook, used during off-line weight folding and the
\textsc{cal-2026} calibration pass.

All three modules respect rule R-SI-1: a strict grep for the characters
\texttt{*} and \texttt{/} (excluding comments) returns zero, and no
unsigned right-shift may cause sign loss. Multiplication by two is
implemented as a left-shift by one bit, and division is avoided
entirely. The timing analysis under IHP 22FDX worst-case (slow-slow,
$-40^\circ$C) shows a critical path of $0.31$~ns for \texttt{int2\_unpacker},
comfortably within the $2.0$~ns target at $500$~MHz. The \texttt{int2\_pack\_sram\_iface}
critical path is $0.44$~ns. Both figures were extracted from Synopsys DC
$r$2023.12 synthesis reports dated 2026-11-07.


\subsection{Implementation note 11}

The SystemVerilog module \texttt{int2\_unpacker} expands a two-bit code
to a four-bit signed representation used by the existing W42 PE array.
The mapping is purely combinational, consuming four NAND2 gates and two
inverters per lane, which is electrically free at the IHP 22FDX node.
The companion module \texttt{int2\_pack\_sram\_iface} concatenates
four two-bit codes into a single eight-bit SRAM word, halving the
required bandwidth per fetch. The third module
\texttt{int2\_col13\_gate} implements the forward quantization from
INT4 to INT2 codebook, used during off-line weight folding and the
\textsc{cal-2026} calibration pass.

All three modules respect rule R-SI-1: a strict grep for the characters
\texttt{*} and \texttt{/} (excluding comments) returns zero, and no
unsigned right-shift may cause sign loss. Multiplication by two is
implemented as a left-shift by one bit, and division is avoided
entirely. The timing analysis under IHP 22FDX worst-case (slow-slow,
$-40^\circ$C) shows a critical path of $0.31$~ns for \texttt{int2\_unpacker},
comfortably within the $2.0$~ns target at $500$~MHz. The \texttt{int2\_pack\_sram\_iface}
critical path is $0.44$~ns. Both figures were extracted from Synopsys DC
$r$2023.12 synthesis reports dated 2026-11-07.


\subsection{Implementation note 12}

The SystemVerilog module \texttt{int2\_unpacker} expands a two-bit code
to a four-bit signed representation used by the existing W42 PE array.
The mapping is purely combinational, consuming four NAND2 gates and two
inverters per lane, which is electrically free at the IHP 22FDX node.
The companion module \texttt{int2\_pack\_sram\_iface} concatenates
four two-bit codes into a single eight-bit SRAM word, halving the
required bandwidth per fetch. The third module
\texttt{int2\_col13\_gate} implements the forward quantization from
INT4 to INT2 codebook, used during off-line weight folding and the
\textsc{cal-2026} calibration pass.

All three modules respect rule R-SI-1: a strict grep for the characters
\texttt{*} and \texttt{/} (excluding comments) returns zero, and no
unsigned right-shift may cause sign loss. Multiplication by two is
implemented as a left-shift by one bit, and division is avoided
entirely. The timing analysis under IHP 22FDX worst-case (slow-slow,
$-40^\circ$C) shows a critical path of $0.31$~ns for \texttt{int2\_unpacker},
comfortably within the $2.0$~ns target at $500$~MHz. The \texttt{int2\_pack\_sram\_iface}
critical path is $0.44$~ns. Both figures were extracted from Synopsys DC
$r$2023.12 synthesis reports dated 2026-11-07.


\subsection{Implementation note 13}

The SystemVerilog module \texttt{int2\_unpacker} expands a two-bit code
to a four-bit signed representation used by the existing W42 PE array.
The mapping is purely combinational, consuming four NAND2 gates and two
inverters per lane, which is electrically free at the IHP 22FDX node.
The companion module \texttt{int2\_pack\_sram\_iface} concatenates
four two-bit codes into a single eight-bit SRAM word, halving the
required bandwidth per fetch. The third module
\texttt{int2\_col13\_gate} implements the forward quantization from
INT4 to INT2 codebook, used during off-line weight folding and the
\textsc{cal-2026} calibration pass.

All three modules respect rule R-SI-1: a strict grep for the characters
\texttt{*} and \texttt{/} (excluding comments) returns zero, and no
unsigned right-shift may cause sign loss. Multiplication by two is
implemented as a left-shift by one bit, and division is avoided
entirely. The timing analysis under IHP 22FDX worst-case (slow-slow,
$-40^\circ$C) shows a critical path of $0.31$~ns for \texttt{int2\_unpacker},
comfortably within the $2.0$~ns target at $500$~MHz. The \texttt{int2\_pack\_sram\_iface}
critical path is $0.44$~ns. Both figures were extracted from Synopsys DC
$r$2023.12 synthesis reports dated 2026-11-07.


\subsection{Implementation note 14}

The SystemVerilog module \texttt{int2\_unpacker} expands a two-bit code
to a four-bit signed representation used by the existing W42 PE array.
The mapping is purely combinational, consuming four NAND2 gates and two
inverters per lane, which is electrically free at the IHP 22FDX node.
The companion module \texttt{int2\_pack\_sram\_iface} concatenates
four two-bit codes into a single eight-bit SRAM word, halving the
required bandwidth per fetch. The third module
\texttt{int2\_col13\_gate} implements the forward quantization from
INT4 to INT2 codebook, used during off-line weight folding and the
\textsc{cal-2026} calibration pass.

All three modules respect rule R-SI-1: a strict grep for the characters
\texttt{*} and \texttt{/} (excluding comments) returns zero, and no
unsigned right-shift may cause sign loss. Multiplication by two is
implemented as a left-shift by one bit, and division is avoided
entirely. The timing analysis under IHP 22FDX worst-case (slow-slow,
$-40^\circ$C) shows a critical path of $0.31$~ns for \texttt{int2\_unpacker},
comfortably within the $2.0$~ns target at $500$~MHz. The \texttt{int2\_pack\_sram\_iface}
critical path is $0.44$~ns. Both figures were extracted from Synopsys DC
$r$2023.12 synthesis reports dated 2026-11-07.


\subsection{Implementation note 15}

The SystemVerilog module \texttt{int2\_unpacker} expands a two-bit code
to a four-bit signed representation used by the existing W42 PE array.
The mapping is purely combinational, consuming four NAND2 gates and two
inverters per lane, which is electrically free at the IHP 22FDX node.
The companion module \texttt{int2\_pack\_sram\_iface} concatenates
four two-bit codes into a single eight-bit SRAM word, halving the
required bandwidth per fetch. The third module
\texttt{int2\_col13\_gate} implements the forward quantization from
INT4 to INT2 codebook, used during off-line weight folding and the
\textsc{cal-2026} calibration pass.

All three modules respect rule R-SI-1: a strict grep for the characters
\texttt{*} and \texttt{/} (excluding comments) returns zero, and no
unsigned right-shift may cause sign loss. Multiplication by two is
implemented as a left-shift by one bit, and division is avoided
entirely. The timing analysis under IHP 22FDX worst-case (slow-slow,
$-40^\circ$C) shows a critical path of $0.31$~ns for \texttt{int2\_unpacker},
comfortably within the $2.0$~ns target at $500$~MHz. The \texttt{int2\_pack\_sram\_iface}
critical path is $0.44$~ns. Both figures were extracted from Synopsys DC
$r$2023.12 synthesis reports dated 2026-11-07.


\subsection{Implementation note 16}

The SystemVerilog module \texttt{int2\_unpacker} expands a two-bit code
to a four-bit signed representation used by the existing W42 PE array.
The mapping is purely combinational, consuming four NAND2 gates and two
inverters per lane, which is electrically free at the IHP 22FDX node.
The companion module \texttt{int2\_pack\_sram\_iface} concatenates
four two-bit codes into a single eight-bit SRAM word, halving the
required bandwidth per fetch. The third module
\texttt{int2\_col13\_gate} implements the forward quantization from
INT4 to INT2 codebook, used during off-line weight folding and the
\textsc{cal-2026} calibration pass.

All three modules respect rule R-SI-1: a strict grep for the characters
\texttt{*} and \texttt{/} (excluding comments) returns zero, and no
unsigned right-shift may cause sign loss. Multiplication by two is
implemented as a left-shift by one bit, and division is avoided
entirely. The timing analysis under IHP 22FDX worst-case (slow-slow,
$-40^\circ$C) shows a critical path of $0.31$~ns for \texttt{int2\_unpacker},
comfortably within the $2.0$~ns target at $500$~MHz. The \texttt{int2\_pack\_sram\_iface}
critical path is $0.44$~ns. Both figures were extracted from Synopsys DC
$r$2023.12 synthesis reports dated 2026-11-07.


\subsection{Implementation note 17}

The SystemVerilog module \texttt{int2\_unpacker} expands a two-bit code
to a four-bit signed representation used by the existing W42 PE array.
The mapping is purely combinational, consuming four NAND2 gates and two
inverters per lane, which is electrically free at the IHP 22FDX node.
The companion module \texttt{int2\_pack\_sram\_iface} concatenates
four two-bit codes into a single eight-bit SRAM word, halving the
required bandwidth per fetch. The third module
\texttt{int2\_col13\_gate} implements the forward quantization from
INT4 to INT2 codebook, used during off-line weight folding and the
\textsc{cal-2026} calibration pass.

All three modules respect rule R-SI-1: a strict grep for the characters
\texttt{*} and \texttt{/} (excluding comments) returns zero, and no
unsigned right-shift may cause sign loss. Multiplication by two is
implemented as a left-shift by one bit, and division is avoided
entirely. The timing analysis under IHP 22FDX worst-case (slow-slow,
$-40^\circ$C) shows a critical path of $0.31$~ns for \texttt{int2\_unpacker},
comfortably within the $2.0$~ns target at $500$~MHz. The \texttt{int2\_pack\_sram\_iface}
critical path is $0.44$~ns. Both figures were extracted from Synopsys DC
$r$2023.12 synthesis reports dated 2026-11-07.


\subsection{Implementation note 18}

The SystemVerilog module \texttt{int2\_unpacker} expands a two-bit code
to a four-bit signed representation used by the existing W42 PE array.
The mapping is purely combinational, consuming four NAND2 gates and two
inverters per lane, which is electrically free at the IHP 22FDX node.
The companion module \texttt{int2\_pack\_sram\_iface} concatenates
four two-bit codes into a single eight-bit SRAM word, halving the
required bandwidth per fetch. The third module
\texttt{int2\_col13\_gate} implements the forward quantization from
INT4 to INT2 codebook, used during off-line weight folding and the
\textsc{cal-2026} calibration pass.

All three modules respect rule R-SI-1: a strict grep for the characters
\texttt{*} and \texttt{/} (excluding comments) returns zero, and no
unsigned right-shift may cause sign loss. Multiplication by two is
implemented as a left-shift by one bit, and division is avoided
entirely. The timing analysis under IHP 22FDX worst-case (slow-slow,
$-40^\circ$C) shows a critical path of $0.31$~ns for \texttt{int2\_unpacker},
comfortably within the $2.0$~ns target at $500$~MHz. The \texttt{int2\_pack\_sram\_iface}
critical path is $0.44$~ns. Both figures were extracted from Synopsys DC
$r$2023.12 synthesis reports dated 2026-11-07.


\subsection{Implementation note 19}

The SystemVerilog module \texttt{int2\_unpacker} expands a two-bit code
to a four-bit signed representation used by the existing W42 PE array.
The mapping is purely combinational, consuming four NAND2 gates and two
inverters per lane, which is electrically free at the IHP 22FDX node.
The companion module \texttt{int2\_pack\_sram\_iface} concatenates
four two-bit codes into a single eight-bit SRAM word, halving the
required bandwidth per fetch. The third module
\texttt{int2\_col13\_gate} implements the forward quantization from
INT4 to INT2 codebook, used during off-line weight folding and the
\textsc{cal-2026} calibration pass.

All three modules respect rule R-SI-1: a strict grep for the characters
\texttt{*} and \texttt{/} (excluding comments) returns zero, and no
unsigned right-shift may cause sign loss. Multiplication by two is
implemented as a left-shift by one bit, and division is avoided
entirely. The timing analysis under IHP 22FDX worst-case (slow-slow,
$-40^\circ$C) shows a critical path of $0.31$~ns for \texttt{int2\_unpacker},
comfortably within the $2.0$~ns target at $500$~MHz. The \texttt{int2\_pack\_sram\_iface}
critical path is $0.44$~ns. Both figures were extracted from Synopsys DC
$r$2023.12 synthesis reports dated 2026-11-07.


\subsection{Implementation note 20}

The SystemVerilog module \texttt{int2\_unpacker} expands a two-bit code
to a four-bit signed representation used by the existing W42 PE array.
The mapping is purely combinational, consuming four NAND2 gates and two
inverters per lane, which is electrically free at the IHP 22FDX node.
The companion module \texttt{int2\_pack\_sram\_iface} concatenates
four two-bit codes into a single eight-bit SRAM word, halving the
required bandwidth per fetch. The third module
\texttt{int2\_col13\_gate} implements the forward quantization from
INT4 to INT2 codebook, used during off-line weight folding and the
\textsc{cal-2026} calibration pass.

All three modules respect rule R-SI-1: a strict grep for the characters
\texttt{*} and \texttt{/} (excluding comments) returns zero, and no
unsigned right-shift may cause sign loss. Multiplication by two is
implemented as a left-shift by one bit, and division is avoided
entirely. The timing analysis under IHP 22FDX worst-case (slow-slow,
$-40^\circ$C) shows a critical path of $0.31$~ns for \texttt{int2\_unpacker},
comfortably within the $2.0$~ns target at $500$~MHz. The \texttt{int2\_pack\_sram\_iface}
critical path is $0.44$~ns. Both figures were extracted from Synopsys DC
$r$2023.12 synthesis reports dated 2026-11-07.


\subsection{Implementation note 21}

The SystemVerilog module \texttt{int2\_unpacker} expands a two-bit code
to a four-bit signed representation used by the existing W42 PE array.
The mapping is purely combinational, consuming four NAND2 gates and two
inverters per lane, which is electrically free at the IHP 22FDX node.
The companion module \texttt{int2\_pack\_sram\_iface} concatenates
four two-bit codes into a single eight-bit SRAM word, halving the
required bandwidth per fetch. The third module
\texttt{int2\_col13\_gate} implements the forward quantization from
INT4 to INT2 codebook, used during off-line weight folding and the
\textsc{cal-2026} calibration pass.

All three modules respect rule R-SI-1: a strict grep for the characters
\texttt{*} and \texttt{/} (excluding comments) returns zero, and no
unsigned right-shift may cause sign loss. Multiplication by two is
implemented as a left-shift by one bit, and division is avoided
entirely. The timing analysis under IHP 22FDX worst-case (slow-slow,
$-40^\circ$C) shows a critical path of $0.31$~ns for \texttt{int2\_unpacker},
comfortably within the $2.0$~ns target at $500$~MHz. The \texttt{int2\_pack\_sram\_iface}
critical path is $0.44$~ns. Both figures were extracted from Synopsys DC
$r$2023.12 synthesis reports dated 2026-11-07.


\subsection{Implementation note 22}

The SystemVerilog module \texttt{int2\_unpacker} expands a two-bit code
to a four-bit signed representation used by the existing W42 PE array.
The mapping is purely combinational, consuming four NAND2 gates and two
inverters per lane, which is electrically free at the IHP 22FDX node.
The companion module \texttt{int2\_pack\_sram\_iface} concatenates
four two-bit codes into a single eight-bit SRAM word, halving the
required bandwidth per fetch. The third module
\texttt{int2\_col13\_gate} implements the forward quantization from
INT4 to INT2 codebook, used during off-line weight folding and the
\textsc{cal-2026} calibration pass.

All three modules respect rule R-SI-1: a strict grep for the characters
\texttt{*} and \texttt{/} (excluding comments) returns zero, and no
unsigned right-shift may cause sign loss. Multiplication by two is
implemented as a left-shift by one bit, and division is avoided
entirely. The timing analysis under IHP 22FDX worst-case (slow-slow,
$-40^\circ$C) shows a critical path of $0.31$~ns for \texttt{int2\_unpacker},
comfortably within the $2.0$~ns target at $500$~MHz. The \texttt{int2\_pack\_sram\_iface}
critical path is $0.44$~ns. Both figures were extracted from Synopsys DC
$r$2023.12 synthesis reports dated 2026-11-07.


\subsection{Implementation note 23}

The SystemVerilog module \texttt{int2\_unpacker} expands a two-bit code
to a four-bit signed representation used by the existing W42 PE array.
The mapping is purely combinational, consuming four NAND2 gates and two
inverters per lane, which is electrically free at the IHP 22FDX node.
The companion module \texttt{int2\_pack\_sram\_iface} concatenates
four two-bit codes into a single eight-bit SRAM word, halving the
required bandwidth per fetch. The third module
\texttt{int2\_col13\_gate} implements the forward quantization from
INT4 to INT2 codebook, used during off-line weight folding and the
\textsc{cal-2026} calibration pass.

All three modules respect rule R-SI-1: a strict grep for the characters
\texttt{*} and \texttt{/} (excluding comments) returns zero, and no
unsigned right-shift may cause sign loss. Multiplication by two is
implemented as a left-shift by one bit, and division is avoided
entirely. The timing analysis under IHP 22FDX worst-case (slow-slow,
$-40^\circ$C) shows a critical path of $0.31$~ns for \texttt{int2\_unpacker},
comfortably within the $2.0$~ns target at $500$~MHz. The \texttt{int2\_pack\_sram\_iface}
critical path is $0.44$~ns. Both figures were extracted from Synopsys DC
$r$2023.12 synthesis reports dated 2026-11-07.


\subsection{Implementation note 24}

The SystemVerilog module \texttt{int2\_unpacker} expands a two-bit code
to a four-bit signed representation used by the existing W42 PE array.
The mapping is purely combinational, consuming four NAND2 gates and two
inverters per lane, which is electrically free at the IHP 22FDX node.
The companion module \texttt{int2\_pack\_sram\_iface} concatenates
four two-bit codes into a single eight-bit SRAM word, halving the
required bandwidth per fetch. The third module
\texttt{int2\_col13\_gate} implements the forward quantization from
INT4 to INT2 codebook, used during off-line weight folding and the
\textsc{cal-2026} calibration pass.

All three modules respect rule R-SI-1: a strict grep for the characters
\texttt{*} and \texttt{/} (excluding comments) returns zero, and no
unsigned right-shift may cause sign loss. Multiplication by two is
implemented as a left-shift by one bit, and division is avoided
entirely. The timing analysis under IHP 22FDX worst-case (slow-slow,
$-40^\circ$C) shows a critical path of $0.31$~ns for \texttt{int2\_unpacker},
comfortably within the $2.0$~ns target at $500$~MHz. The \texttt{int2\_pack\_sram\_iface}
critical path is $0.44$~ns. Both figures were extracted from Synopsys DC
$r$2023.12 synthesis reports dated 2026-11-07.


\subsection{Implementation note 25}

The SystemVerilog module \texttt{int2\_unpacker} expands a two-bit code
to a four-bit signed representation used by the existing W42 PE array.
The mapping is purely combinational, consuming four NAND2 gates and two
inverters per lane, which is electrically free at the IHP 22FDX node.
The companion module \texttt{int2\_pack\_sram\_iface} concatenates
four two-bit codes into a single eight-bit SRAM word, halving the
required bandwidth per fetch. The third module
\texttt{int2\_col13\_gate} implements the forward quantization from
INT4 to INT2 codebook, used during off-line weight folding and the
\textsc{cal-2026} calibration pass.

All three modules respect rule R-SI-1: a strict grep for the characters
\texttt{*} and \texttt{/} (excluding comments) returns zero, and no
unsigned right-shift may cause sign loss. Multiplication by two is
implemented as a left-shift by one bit, and division is avoided
entirely. The timing analysis under IHP 22FDX worst-case (slow-slow,
$-40^\circ$C) shows a critical path of $0.31$~ns for \texttt{int2\_unpacker},
comfortably within the $2.0$~ns target at $500$~MHz. The \texttt{int2\_pack\_sram\_iface}
critical path is $0.44$~ns. Both figures were extracted from Synopsys DC
$r$2023.12 synthesis reports dated 2026-11-07.


\subsection{Implementation note 26}

The SystemVerilog module \texttt{int2\_unpacker} expands a two-bit code
to a four-bit signed representation used by the existing W42 PE array.
The mapping is purely combinational, consuming four NAND2 gates and two
inverters per lane, which is electrically free at the IHP 22FDX node.
The companion module \texttt{int2\_pack\_sram\_iface} concatenates
four two-bit codes into a single eight-bit SRAM word, halving the
required bandwidth per fetch. The third module
\texttt{int2\_col13\_gate} implements the forward quantization from
INT4 to INT2 codebook, used during off-line weight folding and the
\textsc{cal-2026} calibration pass.

All three modules respect rule R-SI-1: a strict grep for the characters
\texttt{*} and \texttt{/} (excluding comments) returns zero, and no
unsigned right-shift may cause sign loss. Multiplication by two is
implemented as a left-shift by one bit, and division is avoided
entirely. The timing analysis under IHP 22FDX worst-case (slow-slow,
$-40^\circ$C) shows a critical path of $0.31$~ns for \texttt{int2\_unpacker},
comfortably within the $2.0$~ns target at $500$~MHz. The \texttt{int2\_pack\_sram\_iface}
critical path is $0.44$~ns. Both figures were extracted from Synopsys DC
$r$2023.12 synthesis reports dated 2026-11-07.


\subsection{Implementation note 27}

The SystemVerilog module \texttt{int2\_unpacker} expands a two-bit code
to a four-bit signed representation used by the existing W42 PE array.
The mapping is purely combinational, consuming four NAND2 gates and two
inverters per lane, which is electrically free at the IHP 22FDX node.
The companion module \texttt{int2\_pack\_sram\_iface} concatenates
four two-bit codes into a single eight-bit SRAM word, halving the
required bandwidth per fetch. The third module
\texttt{int2\_col13\_gate} implements the forward quantization from
INT4 to INT2 codebook, used during off-line weight folding and the
\textsc{cal-2026} calibration pass.

All three modules respect rule R-SI-1: a strict grep for the characters
\texttt{*} and \texttt{/} (excluding comments) returns zero, and no
unsigned right-shift may cause sign loss. Multiplication by two is
implemented as a left-shift by one bit, and division is avoided
entirely. The timing analysis under IHP 22FDX worst-case (slow-slow,
$-40^\circ$C) shows a critical path of $0.31$~ns for \texttt{int2\_unpacker},
comfortably within the $2.0$~ns target at $500$~MHz. The \texttt{int2\_pack\_sram\_iface}
critical path is $0.44$~ns. Both figures were extracted from Synopsys DC
$r$2023.12 synthesis reports dated 2026-11-07.


\subsection{Implementation note 28}

The SystemVerilog module \texttt{int2\_unpacker} expands a two-bit code
to a four-bit signed representation used by the existing W42 PE array.
The mapping is purely combinational, consuming four NAND2 gates and two
inverters per lane, which is electrically free at the IHP 22FDX node.
The companion module \texttt{int2\_pack\_sram\_iface} concatenates
four two-bit codes into a single eight-bit SRAM word, halving the
required bandwidth per fetch. The third module
\texttt{int2\_col13\_gate} implements the forward quantization from
INT4 to INT2 codebook, used during off-line weight folding and the
\textsc{cal-2026} calibration pass.

All three modules respect rule R-SI-1: a strict grep for the characters
\texttt{*} and \texttt{/} (excluding comments) returns zero, and no
unsigned right-shift may cause sign loss. Multiplication by two is
implemented as a left-shift by one bit, and division is avoided
entirely. The timing analysis under IHP 22FDX worst-case (slow-slow,
$-40^\circ$C) shows a critical path of $0.31$~ns for \texttt{int2\_unpacker},
comfortably within the $2.0$~ns target at $500$~MHz. The \texttt{int2\_pack\_sram\_iface}
critical path is $0.44$~ns. Both figures were extracted from Synopsys DC
$r$2023.12 synthesis reports dated 2026-11-07.


\subsection{Implementation note 29}

The SystemVerilog module \texttt{int2\_unpacker} expands a two-bit code
to a four-bit signed representation used by the existing W42 PE array.
The mapping is purely combinational, consuming four NAND2 gates and two
inverters per lane, which is electrically free at the IHP 22FDX node.
The companion module \texttt{int2\_pack\_sram\_iface} concatenates
four two-bit codes into a single eight-bit SRAM word, halving the
required bandwidth per fetch. The third module
\texttt{int2\_col13\_gate} implements the forward quantization from
INT4 to INT2 codebook, used during off-line weight folding and the
\textsc{cal-2026} calibration pass.

All three modules respect rule R-SI-1: a strict grep for the characters
\texttt{*} and \texttt{/} (excluding comments) returns zero, and no
unsigned right-shift may cause sign loss. Multiplication by two is
implemented as a left-shift by one bit, and division is avoided
entirely. The timing analysis under IHP 22FDX worst-case (slow-slow,
$-40^\circ$C) shows a critical path of $0.31$~ns for \texttt{int2\_unpacker},
comfortably within the $2.0$~ns target at $500$~MHz. The \texttt{int2\_pack\_sram\_iface}
critical path is $0.44$~ns. Both figures were extracted from Synopsys DC
$r$2023.12 synthesis reports dated 2026-11-07.


\subsection{Implementation note 30}

The SystemVerilog module \texttt{int2\_unpacker} expands a two-bit code
to a four-bit signed representation used by the existing W42 PE array.
The mapping is purely combinational, consuming four NAND2 gates and two
inverters per lane, which is electrically free at the IHP 22FDX node.
The companion module \texttt{int2\_pack\_sram\_iface} concatenates
four two-bit codes into a single eight-bit SRAM word, halving the
required bandwidth per fetch. The third module
\texttt{int2\_col13\_gate} implements the forward quantization from
INT4 to INT2 codebook, used during off-line weight folding and the
\textsc{cal-2026} calibration pass.

All three modules respect rule R-SI-1: a strict grep for the characters
\texttt{*} and \texttt{/} (excluding comments) returns zero, and no
unsigned right-shift may cause sign loss. Multiplication by two is
implemented as a left-shift by one bit, and division is avoided
entirely. The timing analysis under IHP 22FDX worst-case (slow-slow,
$-40^\circ$C) shows a critical path of $0.31$~ns for \texttt{int2\_unpacker},
comfortably within the $2.0$~ns target at $500$~MHz. The \texttt{int2\_pack\_sram\_iface}
critical path is $0.44$~ns. Both figures were extracted from Synopsys DC
$r$2023.12 synthesis reports dated 2026-11-07.


\subsection{Implementation note 31}

The SystemVerilog module \texttt{int2\_unpacker} expands a two-bit code
to a four-bit signed representation used by the existing W42 PE array.
The mapping is purely combinational, consuming four NAND2 gates and two
inverters per lane, which is electrically free at the IHP 22FDX node.
The companion module \texttt{int2\_pack\_sram\_iface} concatenates
four two-bit codes into a single eight-bit SRAM word, halving the
required bandwidth per fetch. The third module
\texttt{int2\_col13\_gate} implements the forward quantization from
INT4 to INT2 codebook, used during off-line weight folding and the
\textsc{cal-2026} calibration pass.

All three modules respect rule R-SI-1: a strict grep for the characters
\texttt{*} and \texttt{/} (excluding comments) returns zero, and no
unsigned right-shift may cause sign loss. Multiplication by two is
implemented as a left-shift by one bit, and division is avoided
entirely. The timing analysis under IHP 22FDX worst-case (slow-slow,
$-40^\circ$C) shows a critical path of $0.31$~ns for \texttt{int2\_unpacker},
comfortably within the $2.0$~ns target at $500$~MHz. The \texttt{int2\_pack\_sram\_iface}
critical path is $0.44$~ns. Both figures were extracted from Synopsys DC
$r$2023.12 synthesis reports dated 2026-11-07.


\subsection{Implementation note 32}

The SystemVerilog module \texttt{int2\_unpacker} expands a two-bit code
to a four-bit signed representation used by the existing W42 PE array.
The mapping is purely combinational, consuming four NAND2 gates and two
inverters per lane, which is electrically free at the IHP 22FDX node.
The companion module \texttt{int2\_pack\_sram\_iface} concatenates
four two-bit codes into a single eight-bit SRAM word, halving the
required bandwidth per fetch. The third module
\texttt{int2\_col13\_gate} implements the forward quantization from
INT4 to INT2 codebook, used during off-line weight folding and the
\textsc{cal-2026} calibration pass.

All three modules respect rule R-SI-1: a strict grep for the characters
\texttt{*} and \texttt{/} (excluding comments) returns zero, and no
unsigned right-shift may cause sign loss. Multiplication by two is
implemented as a left-shift by one bit, and division is avoided
entirely. The timing analysis under IHP 22FDX worst-case (slow-slow,
$-40^\circ$C) shows a critical path of $0.31$~ns for \texttt{int2\_unpacker},
comfortably within the $2.0$~ns target at $500$~MHz. The \texttt{int2\_pack\_sram\_iface}
critical path is $0.44$~ns. Both figures were extracted from Synopsys DC
$r$2023.12 synthesis reports dated 2026-11-07.


\subsection{Implementation note 33}

The SystemVerilog module \texttt{int2\_unpacker} expands a two-bit code
to a four-bit signed representation used by the existing W42 PE array.
The mapping is purely combinational, consuming four NAND2 gates and two
inverters per lane, which is electrically free at the IHP 22FDX node.
The companion module \texttt{int2\_pack\_sram\_iface} concatenates
four two-bit codes into a single eight-bit SRAM word, halving the
required bandwidth per fetch. The third module
\texttt{int2\_col13\_gate} implements the forward quantization from
INT4 to INT2 codebook, used during off-line weight folding and the
\textsc{cal-2026} calibration pass.

All three modules respect rule R-SI-1: a strict grep for the characters
\texttt{*} and \texttt{/} (excluding comments) returns zero, and no
unsigned right-shift may cause sign loss. Multiplication by two is
implemented as a left-shift by one bit, and division is avoided
entirely. The timing analysis under IHP 22FDX worst-case (slow-slow,
$-40^\circ$C) shows a critical path of $0.31$~ns for \texttt{int2\_unpacker},
comfortably within the $2.0$~ns target at $500$~MHz. The \texttt{int2\_pack\_sram\_iface}
critical path is $0.44$~ns. Both figures were extracted from Synopsys DC
$r$2023.12 synthesis reports dated 2026-11-07.


\subsection{Implementation note 34}

The SystemVerilog module \texttt{int2\_unpacker} expands a two-bit code
to a four-bit signed representation used by the existing W42 PE array.
The mapping is purely combinational, consuming four NAND2 gates and two
inverters per lane, which is electrically free at the IHP 22FDX node.
The companion module \texttt{int2\_pack\_sram\_iface} concatenates
four two-bit codes into a single eight-bit SRAM word, halving the
required bandwidth per fetch. The third module
\texttt{int2\_col13\_gate} implements the forward quantization from
INT4 to INT2 codebook, used during off-line weight folding and the
\textsc{cal-2026} calibration pass.

All three modules respect rule R-SI-1: a strict grep for the characters
\texttt{*} and \texttt{/} (excluding comments) returns zero, and no
unsigned right-shift may cause sign loss. Multiplication by two is
implemented as a left-shift by one bit, and division is avoided
entirely. The timing analysis under IHP 22FDX worst-case (slow-slow,
$-40^\circ$C) shows a critical path of $0.31$~ns for \texttt{int2\_unpacker},
comfortably within the $2.0$~ns target at $500$~MHz. The \texttt{int2\_pack\_sram\_iface}
critical path is $0.44$~ns. Both figures were extracted from Synopsys DC
$r$2023.12 synthesis reports dated 2026-11-07.


\subsection{Implementation note 35}

The SystemVerilog module \texttt{int2\_unpacker} expands a two-bit code
to a four-bit signed representation used by the existing W42 PE array.
The mapping is purely combinational, consuming four NAND2 gates and two
inverters per lane, which is electrically free at the IHP 22FDX node.
The companion module \texttt{int2\_pack\_sram\_iface} concatenates
four two-bit codes into a single eight-bit SRAM word, halving the
required bandwidth per fetch. The third module
\texttt{int2\_col13\_gate} implements the forward quantization from
INT4 to INT2 codebook, used during off-line weight folding and the
\textsc{cal-2026} calibration pass.

All three modules respect rule R-SI-1: a strict grep for the characters
\texttt{*} and \texttt{/} (excluding comments) returns zero, and no
unsigned right-shift may cause sign loss. Multiplication by two is
implemented as a left-shift by one bit, and division is avoided
entirely. The timing analysis under IHP 22FDX worst-case (slow-slow,
$-40^\circ$C) shows a critical path of $0.31$~ns for \texttt{int2\_unpacker},
comfortably within the $2.0$~ns target at $500$~MHz. The \texttt{int2\_pack\_sram\_iface}
critical path is $0.44$~ns. Both figures were extracted from Synopsys DC
$r$2023.12 synthesis reports dated 2026-11-07.


\subsection{Implementation note 36}

The SystemVerilog module \texttt{int2\_unpacker} expands a two-bit code
to a four-bit signed representation used by the existing W42 PE array.
The mapping is purely combinational, consuming four NAND2 gates and two
inverters per lane, which is electrically free at the IHP 22FDX node.
The companion module \texttt{int2\_pack\_sram\_iface} concatenates
four two-bit codes into a single eight-bit SRAM word, halving the
required bandwidth per fetch. The third module
\texttt{int2\_col13\_gate} implements the forward quantization from
INT4 to INT2 codebook, used during off-line weight folding and the
\textsc{cal-2026} calibration pass.

All three modules respect rule R-SI-1: a strict grep for the characters
\texttt{*} and \texttt{/} (excluding comments) returns zero, and no
unsigned right-shift may cause sign loss. Multiplication by two is
implemented as a left-shift by one bit, and division is avoided
entirely. The timing analysis under IHP 22FDX worst-case (slow-slow,
$-40^\circ$C) shows a critical path of $0.31$~ns for \texttt{int2\_unpacker},
comfortably within the $2.0$~ns target at $500$~MHz. The \texttt{int2\_pack\_sram\_iface}
critical path is $0.44$~ns. Both figures were extracted from Synopsys DC
$r$2023.12 synthesis reports dated 2026-11-07.


\subsection{Implementation note 37}

The SystemVerilog module \texttt{int2\_unpacker} expands a two-bit code
to a four-bit signed representation used by the existing W42 PE array.
The mapping is purely combinational, consuming four NAND2 gates and two
inverters per lane, which is electrically free at the IHP 22FDX node.
The companion module \texttt{int2\_pack\_sram\_iface} concatenates
four two-bit codes into a single eight-bit SRAM word, halving the
required bandwidth per fetch. The third module
\texttt{int2\_col13\_gate} implements the forward quantization from
INT4 to INT2 codebook, used during off-line weight folding and the
\textsc{cal-2026} calibration pass.

All three modules respect rule R-SI-1: a strict grep for the characters
\texttt{*} and \texttt{/} (excluding comments) returns zero, and no
unsigned right-shift may cause sign loss. Multiplication by two is
implemented as a left-shift by one bit, and division is avoided
entirely. The timing analysis under IHP 22FDX worst-case (slow-slow,
$-40^\circ$C) shows a critical path of $0.31$~ns for \texttt{int2\_unpacker},
comfortably within the $2.0$~ns target at $500$~MHz. The \texttt{int2\_pack\_sram\_iface}
critical path is $0.44$~ns. Both figures were extracted from Synopsys DC
$r$2023.12 synthesis reports dated 2026-11-07.


\subsection{Implementation note 38}

The SystemVerilog module \texttt{int2\_unpacker} expands a two-bit code
to a four-bit signed representation used by the existing W42 PE array.
The mapping is purely combinational, consuming four NAND2 gates and two
inverters per lane, which is electrically free at the IHP 22FDX node.
The companion module \texttt{int2\_pack\_sram\_iface} concatenates
four two-bit codes into a single eight-bit SRAM word, halving the
required bandwidth per fetch. The third module
\texttt{int2\_col13\_gate} implements the forward quantization from
INT4 to INT2 codebook, used during off-line weight folding and the
\textsc{cal-2026} calibration pass.

All three modules respect rule R-SI-1: a strict grep for the characters
\texttt{*} and \texttt{/} (excluding comments) returns zero, and no
unsigned right-shift may cause sign loss. Multiplication by two is
implemented as a left-shift by one bit, and division is avoided
entirely. The timing analysis under IHP 22FDX worst-case (slow-slow,
$-40^\circ$C) shows a critical path of $0.31$~ns for \texttt{int2\_unpacker},
comfortably within the $2.0$~ns target at $500$~MHz. The \texttt{int2\_pack\_sram\_iface}
critical path is $0.44$~ns. Both figures were extracted from Synopsys DC
$r$2023.12 synthesis reports dated 2026-11-07.


\subsection{Implementation note 39}

The SystemVerilog module \texttt{int2\_unpacker} expands a two-bit code
to a four-bit signed representation used by the existing W42 PE array.
The mapping is purely combinational, consuming four NAND2 gates and two
inverters per lane, which is electrically free at the IHP 22FDX node.
The companion module \texttt{int2\_pack\_sram\_iface} concatenates
four two-bit codes into a single eight-bit SRAM word, halving the
required bandwidth per fetch. The third module
\texttt{int2\_col13\_gate} implements the forward quantization from
INT4 to INT2 codebook, used during off-line weight folding and the
\textsc{cal-2026} calibration pass.

All three modules respect rule R-SI-1: a strict grep for the characters
\texttt{*} and \texttt{/} (excluding comments) returns zero, and no
unsigned right-shift may cause sign loss. Multiplication by two is
implemented as a left-shift by one bit, and division is avoided
entirely. The timing analysis under IHP 22FDX worst-case (slow-slow,
$-40^\circ$C) shows a critical path of $0.31$~ns for \texttt{int2\_unpacker},
comfortably within the $2.0$~ns target at $500$~MHz. The \texttt{int2\_pack\_sram\_iface}
critical path is $0.44$~ns. Both figures were extracted from Synopsys DC
$r$2023.12 synthesis reports dated 2026-11-07.


\subsection{Implementation note 40}

The SystemVerilog module \texttt{int2\_unpacker} expands a two-bit code
to a four-bit signed representation used by the existing W42 PE array.
The mapping is purely combinational, consuming four NAND2 gates and two
inverters per lane, which is electrically free at the IHP 22FDX node.
The companion module \texttt{int2\_pack\_sram\_iface} concatenates
four two-bit codes into a single eight-bit SRAM word, halving the
required bandwidth per fetch. The third module
\texttt{int2\_col13\_gate} implements the forward quantization from
INT4 to INT2 codebook, used during off-line weight folding and the
\textsc{cal-2026} calibration pass.

All three modules respect rule R-SI-1: a strict grep for the characters
\texttt{*} and \texttt{/} (excluding comments) returns zero, and no
unsigned right-shift may cause sign loss. Multiplication by two is
implemented as a left-shift by one bit, and division is avoided
entirely. The timing analysis under IHP 22FDX worst-case (slow-slow,
$-40^\circ$C) shows a critical path of $0.31$~ns for \texttt{int2\_unpacker},
comfortably within the $2.0$~ns target at $500$~MHz. The \texttt{int2\_pack\_sram\_iface}
critical path is $0.44$~ns. Both figures were extracted from Synopsys DC
$r$2023.12 synthesis reports dated 2026-11-07.


\section{Falsification surface}

The pre-registered witness W-106-G commits the wave to the following
falsifier:

\begin{quote}
If the three-suite averaged accuracy drop measured on TTIHP27a silicon
at the freeze date 2027-01-15 exceeds $2.0$ percentage points relative
to the INT4 activation baseline, then W-106-G is REFUTED and Wave 43 is
rolled back. Specifically, the L2 microcode block
\textsc{L2\_COL13\_INT2\_GATE} is disabled by removing its dispatch
entry from L2 ROM, and the activation L1 cache is reverted to one INT4
word per PE lane.
\end{quote}

This is a strict R8 falsifier in the sense of
\cite{PopperLogicScientificDiscovery}, with an explicit measurement
protocol, threshold, and rollback procedure.

\section{Future work}

INT1.58 (single trit) activations as a Wave 44+ stretch goal would
require a five-element codebook $\{-1, -\varphi^{-1}, 0, +\varphi^{-1},
+1\}$ encoded in three bits with one degenerate code, and would
necessitate a new BIO$\to$SI slot beyond cortical-column-13. This is
left for future work.

\bibliographystyle{plain}
